\documentclass[11pt]{article}

\usepackage[letterpaper,margin=0.82in,headheight=15pt]{geometry}
\usepackage[T1]{fontenc}
\usepackage{lmodern}
\usepackage{microtype}
\usepackage{amsmath,amssymb,amsthm,bm,mathtools}
\usepackage{stmaryrd}
\usepackage{booktabs,longtable,tabularx,array,multirow}
\usepackage{enumitem}
\usepackage{xcolor}
\usepackage{fancyhdr}
\usepackage{graphicx}
\usepackage{tikz-cd}
\usepackage{hyperref}
\usepackage[nameinlink,noabbrev]{cleveref}

\definecolor{navy}{HTML}{17324D}
\definecolor{teal}{HTML}{147D92}
\definecolor{orange}{HTML}{D47A27}
\definecolor{softblue}{HTML}{EEF4F7}
\definecolor{softgray}{HTML}{F3F5F6}
\definecolor{darkgray}{HTML}{46545E}
\definecolor{softorange}{HTML}{FFF4E8}
\definecolor{softgreen}{HTML}{EDF7F0}
\definecolor{softred}{HTML}{FCEEEE}

\hypersetup{
  colorlinks=true,
  linkcolor=navy,
  citecolor=teal,
  urlcolor=teal,
  pdftitle={Volume VIII: Symmetry-Adapted Tensor Networks and the Assumption-Conditioned Prediction Compiler},
  pdfauthor={DeuteronWigner project}
}

\pagestyle{fancy}
\fancyhf{}
\fancyhead[L]{\small\color{darkgray}DeuteronWigner formalism - Volume VIII}
\fancyhead[R]{\small\color{darkgray}Tensor networks and prediction compiler}
\fancyfoot[L]{\small\color{darkgray}30 July 2026}
\fancyfoot[R]{\small\color{darkgray}\thepage}
\renewcommand{\headrulewidth}{0.4pt}

\setlist[itemize]{leftmargin=1.5em,itemsep=2pt,topsep=3pt}
\setlist[enumerate]{leftmargin=1.7em,itemsep=2pt,topsep=3pt}
\setlength{\parindent}{0pt}
\setlength{\parskip}{5pt}
\setlength{\emergencystretch}{2em}
\renewcommand{\arraystretch}{1.16}
\allowdisplaybreaks

\newtheorem{definition}{Definition}[section]
\newtheorem{axiom}[definition]{Axiom}
\newtheorem{principle}[definition]{Design principle}
\newtheorem{requirement}[definition]{Requirement}
\newtheorem{proposition}[definition]{Proposition}
\newtheorem{theorem}[definition]{Theorem}
\newtheorem{lemma}[definition]{Lemma}
\newtheorem{corollary}[definition]{Corollary}
\newtheorem{remark}[definition]{Remark}
\newtheorem{decision}[definition]{Model decision}

\crefname{definition}{definition}{definitions}
\crefname{axiom}{axiom}{axioms}
\crefname{principle}{principle}{principles}
\crefname{requirement}{requirement}{requirements}
\crefname{proposition}{proposition}{propositions}
\crefname{theorem}{theorem}{theorems}
\crefname{lemma}{lemma}{lemmas}
\crefname{corollary}{corollary}{corollaries}
\crefname{remark}{remark}{remarks}
\crefname{decision}{model decision}{model decisions}

\newcommand{\kT}{\bm{k}_{T}}
\newcommand{\pT}{\bm{p}_{T}}
\newcommand{\qT}{\bm{q}_{T}}
\newcommand{\DT}{\bm{\Delta}_{T}}
\newcommand{\bD}{\bm{b}_{\Delta}}
\newcommand{\bM}{\bm{b}_{\mathrm{TMD}}}
\newcommand{\dd}{\mathrm{d}}
\newcommand{\Tr}{\operatorname{Tr}}
\newcommand{\tr}{\operatorname{tr}}
\newcommand{\diag}{\operatorname{diag}}
\newcommand{\rank}{\operatorname{rank}}
\newcommand{\Span}{\operatorname{span}}
\newcommand{\Ker}{\operatorname{Ker}}
\newcommand{\Ran}{\operatorname{Ran}}
\newcommand{\Cov}{\operatorname{Cov}}
\newcommand{\Var}{\operatorname{Var}}
\newcommand{\Id}{\operatorname{Id}}
\newcommand{\Hom}{\operatorname{Hom}}
\newcommand{\End}{\operatorname{End}}
\newcommand{\Obj}{\operatorname{Obj}}
\newcommand{\Hil}{\mathcal H}
\newcommand{\LF}{\ensuremath{\mathrm{LF}}}
\newcommand{\BLFQ}{\ensuremath{\mathrm{BLFQ}}}
\newcommand{\TMD}{\ensuremath{\mathrm{TMD}}}
\newcommand{\GTMD}{\ensuremath{\mathrm{GTMD}}}
\newcommand{\GPD}{\ensuremath{\mathrm{GPD}}}
\newcommand{\QCD}{\ensuremath{\mathrm{QCD}}}
\newcommand{\EFT}{\ensuremath{\mathrm{EFT}}}
\newcommand{\TTN}{\ensuremath{\mathrm{TTN}}}
\newcommand{\MPS}{\ensuremath{\mathrm{MPS}}}
\newcommand{\PEPS}{\ensuremath{\mathrm{PEPS}}}
\newcommand{\TN}{\ensuremath{\mathrm{TN}}}
\newcommand{\Amp}{\mathbf{Amp}}
\newcommand{\Dens}{\mathbf{Dens}}
\newcommand{\Match}{\mathbf{Match}}
\newcommand{\Red}{\mathbf{Red}}
\newcommand{\Proc}{\mathbf{Proc}}
\newcommand{\Infer}{\mathbf{Infer}}
\newcommand{\Assump}{\mathbf{Assump}}
\newcommand{\PlanCat}{\mathbf{Plan}}
\newcommand{\Prov}{\mathbf{Prov}}
\newcommand{\cO}{\mathcal O}
\newcommand{\cM}{\mathcal M}
\newcommand{\cR}{\mathcal R}
\newcommand{\cS}{\mathcal S}
\newcommand{\cV}{\mathcal V}
\newcommand{\cW}{\mathcal W}
\newcommand{\cE}{\mathcal E}
\newcommand{\cG}{\mathcal G}
\newcommand{\cK}{\mathcal K}
\newcommand{\cQ}{\mathcal Q}
\newcommand{\cP}{\mathcal P}
\newcommand{\cC}{\mathcal C}
\newcommand{\cZ}{\mathcal Z}
\newcommand{\cD}{\mathcal D}
\newcommand{\cN}{\mathcal N}
\newcommand{\cU}{\mathcal U}
\newcommand{\cF}{\mathcal F}
\newcommand{\cA}{\mathcal A}
\newcommand{\cB}{\mathcal B}
\newcommand{\cL}{\mathcal L}
\newcommand{\cJ}{\mathcal J}
\newcommand{\cI}{\mathcal I}
\newcommand{\cH}{\mathcal H}
\newcommand{\cT}{\mathcal T}
\newcommand{\cX}{\mathcal X}
\newcommand{\cY}{\mathcal Y}
\newcommand{\cRho}{\mathcal R}
\newcommand{\eps}{\varepsilon}
\newcommand{\bbone}{\bm 1}
\newcommand{\abs}[1]{\left|#1\right|}
\newcommand{\norm}[1]{\left\lVert#1\right\rVert}
\newcommand{\inner}[2]{\left\langle #1\middle|#2\right\rangle}
\newcommand{\avg}[1]{\left\langle #1\right\rangle}
\newcommand{\phys}{\mathrm{phys}}
\newcommand{\eff}{\mathrm{eff}}
\newcommand{\bare}{\mathrm{bare}}
\newcommand{\ind}{\mathrm{ind}}
\newcommand{\ct}{\mathrm{ct}}
\newcommand{\reg}{\mathrm{reg}}
\newcommand{\trunc}{\mathrm{trunc}}
\newcommand{\ren}{\mathrm{ren}}
\newcommand{\num}{\mathrm{num}}
\newcommand{\proc}{\mathrm{proc}}
\newcommand{\stat}{\mathrm{stat}}
\newcommand{\calib}{\mathrm{cal}}
\newcommand{\hold}{\mathrm{hold}}
\newcommand{\pilot}{\mathrm{pilot}}
\newcommand{\UV}{\mathrm{UV}}
\newcommand{\IR}{\mathrm{IR}}
\newcommand{\COM}{\mathrm{CM}}
\newcommand{\Lz}{L_z}
\newcommand{\Jz}{J^z}
\newcommand{\chiTN}{\chi}
\newcommand{\Aset}{\mathfrak A}
\newcommand{\Pplan}{\mathfrak P}
\newcommand{\Nstate}{\mathcal N_{\mathrm{state}}}
\newcommand{\Nop}{\mathcal N_{\mathrm{op}}}
\newcommand{\Kprov}{K_{\mathrm{prov}}}
\newcommand{\Compiler}{\texttt{PredictionCompiler}}
\newcommand{\AssumptionBundle}{\texttt{AssumptionBundle}}
\newcommand{\PredictionPlan}{\texttt{PredictionPlan}}
\newcommand{\TensorEdgeType}{\texttt{TensorEdgeType}}
\newcommand{\SymTensor}{\texttt{SymmetryBlockTensor}}
\newcommand{\FusionTree}{\texttt{FusionTree}}
\newcommand{\StateTN}{\texttt{StateTensorNetwork}}
\newcommand{\OperatorTN}{\texttt{OperatorTensorNetwork}}
\newcommand{\ProvComplex}{\texttt{Provenance2Complex}}
\newcommand{\ContractionPlan}{\texttt{ContractionPlan}}

\newcolumntype{Y}{>{\raggedright\arraybackslash}X}
\newcolumntype{L}[1]{>{\raggedright\arraybackslash}p{#1}}

\title{\vspace{-1.0cm}
{\color{navy}\bfseries Volume VIII: Symmetry-Adapted Tensor Networks and the\\[2mm]
Assumption-Conditioned Prediction Compiler}\\[4mm]
\large Equivariant tensor blocks, topology-aware contraction graphs, variational light-front states,\\
compositional provenance, and conditional GTMD/TMD prediction plans}
\author{DeuteronWigner project}
\date{30 July 2026}

\begin{document}
\maketitle

\begin{center}
\begin{minipage}{0.94\textwidth}
\colorbox{softblue}{\parbox{0.96\textwidth}{
\textbf{Status and purpose.}
Volumes~0--VII define the algebraic and geometric architecture, regulated
light-front state spaces, common nucleon GTMD operators, dynamical Wilson-line
phases, matched spin-1 nuclear structure, regulator-to-QCD matching,
evolution, inference, and the selected microscopic Hamiltonian route.  C7/H0
has established an exact symmetry-complete finite basis and Hamiltonian-term
spine, while C8/H1 is the active implementation package for the first
interacting valence eigenstate, renormalization-flow benchmark, and
symmetry-adapted tree-tensor-network solver.  The present volume formalizes the
missing compositional layer.  It defines how a typed bundle of scientific
assumptions compiles into a symmetry-preserving tensor network, an executable
operator-contraction plan, a provenance two-complex, and a conditional
prediction ensemble.  The network is not permitted to invent dynamics, erase
operator identities, or turn mutually exclusive model branches into additive
contributions.  Its role is to make every choice, approximation, bond
truncation, contraction, and reduction explicit so that one microscopic state
can generate correlated quark, antiquark, gluon, nucleon, deuteron, and process
predictions without introducing an independent coefficient for each named
TMD.
}}
\end{minipage}
\end{center}

\begin{abstract}
The original DeuteronWigner objective is an assumption-explicit generative
model in which a regulated light-front Hamiltonian produces correlated proton,
neutron, and deuteron states, common quark, antiquark, and gluon GTMDs, and all
TMD, GPD, PDF, form-factor, Wilson-line, and process reductions.  This volume
makes the tensor-network and compilation content of that objective precise.
Three structures are kept distinct: a symmetry-adapted state tensor network, an
operator and contraction network, and a provenance two-complex that records
model choices, replacements, exclusions, matching identities, and
no-double-counting cells.  Every tensor edge is a typed direct sum of
representation and multiplicity spaces carrying Fock sector, parton species,
color, flavor, helicity, orbital angular momentum, momentum basis, regulator,
and scheme data.  Equivariant tensors factor into structural Clebsch--Gordan
objects and reduced degeneracy tensors, so forbidden symmetry blocks are absent
rather than numerically suppressed.  A finite light-front state admits an exact
tree-tensor representation when every virtual bond contains the full symmetry-
resolved Schmidt support; lower bond dimensions define a controlled variational
truncation whose energy, state fidelity, current matrix elements, OAM content,
and GTMD reductions must converge separately.  A frozen
\texttt{AssumptionBundle} declares the Hamiltonian route, Fock content,
regulators, induced interactions, tensor topology, bond policy, Wilson and
nuclear options, matching scheme, process status, and inference role.  A
fail-closed compiler computes dependency closure, rejects inconsistent or
unsupported combinations, instantiates the physical tensor and operator
networks, attaches provenance cells, chooses contraction and solver plans, and
returns a versioned \texttt{PredictionPlan}.  Conditional predictions are
written as distributions given the assumption bundle; within-branch posterior
uncertainty is separated from between-branch structural uncertainty.  The
document defines network gauge fixing, variational optimization, adaptive bond
allocation, differentiable contractions, common-parent GTMD/TMD evaluation,
open-system interfaces to Wilson, nuclear, matching, evolution, and inference
layers, formal requirements, analytic benchmarks, negative tests, software
schemas, and staged acceptance gates.  The resulting compiler is the intended
execution layer for the geometric, topological, and tensor-network model
proposed at the beginning of the project.
\end{abstract}

\tableofcontents

\section{Scientific purpose and boundary}
\label{sec:purpose}

\subsection{The missing execution layer}

The previous volumes define the objects that a predictive model must possess,
but a user-facing generative system requires an additional map:
\begin{equation}
 \Aset
 \xrightarrow{\;\Compiler\;}
 \Pplan
 \xrightarrow{\;\mathrm{execute}\;}
 \left\{
 |\Psi_N\rangle,
 |\Psi_D\rangle,
 W_{a/N},W_{a/D},
 F_A,\cO_{\mathrm{process}}
 \right\}.
\label{eq:compiler-chain}
\end{equation}
Here \(\Aset\) is a complete typed assumption bundle and \(\Pplan\) is an
executable prediction plan.  The first arrow is not numerical fitting.  It is a
compilation step that determines whether the requested model is mathematically
and physically admissible, which objects must be built, which alternatives are
exclusive, which matching maps are required, and which predictions are
actually licensed by the declared readiness state.

C7/H0 provides the exact finite basis, color multiplicities, fermionic
antisymmetry, free invariant-mass operator, and reduced sector-changing term
interface.  C8/H1 is designed to add an interacting valence Hamiltonian,
renormalization flow, Hamiltonian-consistent currents, state tracking, exact
and Krylov solvers, and a symmetry-adapted tree tensor network.  This volume
does not assume that C8 has succeeded.  It defines the formal contract that C8
and all later microscopic packages must satisfy.

\begin{definition}[Assumption-conditioned prediction system]
An assumption-conditioned prediction system is a tuple
\begin{equation}
 \mathfrak S=
 \left(
 \AssumptionBundle,
 \Compiler,
 \PredictionPlan,
 \Nstate,
 \Nop,
 \Kprov,
 \cE
 \right),
\label{eq:system-tuple}
\end{equation}
where \(\cE\) is an execution semantics mapping every accepted plan to a
versioned result envelope.  The system is \emph{traceable} when every numerical
entry in the result envelope has a path to its state amplitudes, operators,
parameters, matching maps, assumptions, and approximation manifests.
\end{definition}

\subsection{The original target, stated operationally}

The final model should support a request such as:
\begin{quote}
Use the RSA--BLFQ Hamiltonian with the declared \(qqq\oplus qqqg\oplus
qqqq\bar q\) Fock tower, induced confinement refitted at every resolution,
Wilson order two with a specified rapidity regulator, the matched nuclear-EFT
deuteron route, a stated TMD scheme and accuracy tuple, and a frozen inference
posterior; then predict the full spin-1 TMD family and selected SIDIS or
Drell--Yan observables.
\end{quote}
The system must not interpret this request as permission to combine every
available model component.  It must prove compatibility, construct the
corresponding state and operator networks, reject unsupported process maps,
and preserve branch identity in the output.

\subsection{Explicit nonclaims}

This volume does not claim:
\begin{itemize}
 \item that low tensor-network rank is guaranteed for QCD;
 \item that a tensor topology determines the Hamiltonian or its counterterms;
 \item that network gauge freedom is a physical gauge symmetry;
 \item that provenance homology supplies a physical amplitude;
 \item that all model branches should be Bayesian-averaged;
 \item that the C7 basis or C8 valence benchmark is a converged physical
 nucleon;
 \item that a successful tensor contraction supplies LF-to-QCD matching,
 TMD evolution, nuclear dynamics, or process factorization;
 \item that one scalar bond dimension can measure every truncation error.
\end{itemize}

\begin{principle}[Structure before compression]
Symmetry, operator identity, color, statistics, Fock provenance, and map class
are exact structural data.  They are imposed before any low-rank approximation.
Bond truncation may reduce degeneracy spaces; it may not delete or merge
required representation sectors silently.
\end{principle}

\section{Three networks that must remain distinct}
\label{sec:three-networks}

\subsection{State tensor network}

The state network \(\Nstate\) represents a vector in a declared regulated
physical Hilbert space:
\begin{equation}
 \Nstate(\Aset,\theta,\chi)
 \longmapsto
 |\Psi(\Aset,\theta,\chi)\rangle
 \in \Hil_{\phys}^{(r)}.
\label{eq:state-network}
\end{equation}
Its nodes are reduced amplitudes or isometric fusion maps; its edges carry
physical and virtual representation labels; and its topology and bond spaces
are approximation data.  The state network is the object optimized by the
Rayleigh quotient.

\subsection{Operator and contraction network}

The operator network \(\Nop\) represents Hamiltonian terms, currents, GTMD
insertions, Wilson interactions, nuclear maps, matching kernels, or process
operators.  A contraction plan joins \(\Nop\) to one or more state networks:
\begin{equation}
 \cC_{\Pplan}
 \left[
 \Nstate^{\dagger},\Nop,\Nstate
 \right]
 =
 \langle\Psi'|\widehat{\cO}|\Psi\rangle.
\label{eq:operator-contraction}
\end{equation}
The contraction network may be a tree, a directed acyclic graph, or a general
hypergraph even when the state ansatz itself is a tree.  Contraction order is a
numerical optimization and must not change the physical operator.

\subsection{Assumption and provenance network}

The provenance complex \(\Kprov\) contains typed objects, physical or
descriptive maps, and declared two-cells relating alternative paths.  It
records why a tensor, term, or prediction is present.  It is not contracted as
a quantum amplitude.  Its coefficients are identities, equivalence labels,
subtractions, replacements, exclusions, and provenance records rather than
wave-function amplitudes.

\begin{requirement}[No network conflation]
A state tensor, an operator tensor, and a provenance cell may refer to one
another by immutable identifiers, but their numerical values and composition
laws are different.  The software must reject any attempt to use a provenance
edge as an amplitude, a contraction path as a Wilson path, or a virtual tensor
gauge transformation as a QCD gauge transformation.
\end{requirement}

\subsection{Commuting architecture}

The intended relations are summarized by
\begin{equation}
\begin{tikzcd}[column sep=large,row sep=large]
 \Aset \arrow[r,"\Compiler"] \arrow[d,"\mathrm{provenance}"']
 & \Pplan \arrow[r,"\mathrm{instantiate}"] \arrow[d]
 & (\Nstate,\Nop,\cC) \arrow[d,"\mathrm{contract}"] \\
 \Kprov \arrow[r,"\mathrm{certificate}"']
 & \mathrm{Admissibility\ record} \arrow[r]
 & \mathrm{Result\ envelope}
\end{tikzcd}
\label{eq:three-network-diagram}
\end{equation}
The bottom path certifies the upper path; it does not evaluate it.

\section{The assumption bundle}
\label{sec:assumption-bundle}

\subsection{Six classes of declared choice}

A complete assumption bundle is decomposed as
\begin{equation}
 \Aset=
 \left(
 \Aset_{\phys},
 \Aset_{\reg},
 \Aset_{\mathrm{alg}},
 \Aset_{\num},
 \Aset_{\proc},
 \Aset_{\stat}
 \right).
\label{eq:assumption-decomposition}
\end{equation}
The classes have different scientific meanings:
\begin{longtable}{L{0.17\textwidth}L{0.34\textwidth}L{0.39\textwidth}}
\toprule
Class & Examples & Interpretation\\
\midrule
\endhead
\(\Aset_{\phys}\) & Hamiltonian terms, Fock sectors, current basis, chiral and nuclear sectors, Wilson order & Physical model content or controlled EFT truncation.\\
\(\Aset_{\reg}\) & \(K,N_{\max},b,\lambda,\Lambda_{\UV},\Lambda_{\IR}\), endpoint and rapidity regulators & Resolution and regularization identity.\\
\(\Aset_{\mathrm{alg}}\) & color fusion tree, permutation representation, tensor topology, operator basis, map class & Exact algebraic representation and compositional structure.\\
\(\Aset_{\num}\) & eigensolver, bond policy, contraction path, quadrature, tolerances, precision & Numerical approximation and resource choices.\\
\(\Aset_{\proc}\) & link map, color weights, hard channel, soft prescription, factorization and Glauber status & Process qualification and observable assembly.\\
\(\Aset_{\stat}\) & calibration split, priors, discrepancy basis, posterior member, model-comparison rule & Inference and prediction semantics.\\
\bottomrule
\end{longtable}

An exact convention or symmetry is not a selectable model branch.  For
example, the light-front normalization and the anti-fundamental antiquark
color action are fixed contracts.  By contrast, inclusion of an explicit
\(qqqg\) sector versus its induced effective operator is a physical alternative
subject to a no-double-counting relation.

\subsection{Immutable identity}

An \AssumptionBundle\ is frozen after validation.  Its identity is a content
hash of canonical serialized fields and referenced source hashes:
\begin{equation}
 \mathrm{id}(\Aset)
 =
 \mathrm{SHA256}
 \left[
 \mathrm{canonical}(\Aset),
 \{h_{\mathrm{source}}\}
 \right].
\label{eq:assumption-hash}
\end{equation}
Changing one regulator, Fock sector, Wilson order, tensor topology, or
calibration role creates a new assumption identity.  Mutable aliases are not
accepted as scientific provenance.

\subsection{Choice axes and branch semantics}

Every selectable field belongs to one of the following kinds:
\begin{enumerate}
 \item \textbf{Required singleton}: one exact convention or mandatory object.
 \item \textbf{Exclusive branch}: exactly one of several physical descriptions.
 \item \textbf{Optional additive mechanism}: included only when its source and
 target types permit addition and no replacement relation forbids it.
 \item \textbf{Nested truncation}: an ordered family such as Fock content,
 basis size, Wilson order, or bond dimension.
 \item \textbf{Continuous parameter}: a Hamiltonian, matching, or inference
 coefficient owned by a declared operator.
 \item \textbf{Diagnostic variation}: a noncentral alternative used for
 convergence or sensitivity rather than composition.
\end{enumerate}
A branch choice is not a numerical coefficient.  Setting the coefficient of
an excluded mechanism to zero does not make the branch relation disappear;
the provenance and operator basis remain different.

\subsection{Refinement order}

A partial order \(\Aset\preceq\Aset'\) is permitted only when
\(\Aset'\) refines \(\Aset\) along declared nested axes without changing an
exclusive physical branch:
\begin{equation}
 \Aset\preceq\Aset'
 \Longrightarrow
 \begin{cases}
 \cF_{\Aset}\subseteq\cF_{\Aset'},\\
 r_{\Aset}\preceq r_{\Aset'},\\
 \chi_e(\Aset)\leq\chi_e(\Aset')\quad\forall e,\\
 \text{same branch identities and renormalization conditions}.
 \end{cases}
\label{eq:assumption-refinement}
\end{equation}
Changing from induced confinement to zero confinement is not a refinement; it
creates a sibling branch.

\section{Admissibility logic and compilation}
\label{sec:compiler}

\subsection{Typed predicates}

Let \(\cR=\{R_1,\ldots,R_m\}\) be a finite set of admissibility rules.  Each
rule is a typed predicate
\begin{equation}
 R_i:\Aset\times\cT_{\mathrm{repo}}
 \longrightarrow
 \{\mathrm{pass},\mathrm{fail},\mathrm{unresolved}\},
\label{eq:typed-rule}
\end{equation}
where \(\cT_{\mathrm{repo}}\) denotes the available typed objects and
readiness records.  Rules may encode:
\begin{itemize}
 \item dependencies, such as \(qqqg\Rightarrow\) adjoint color and gluon basis;
 \item exclusions, such as explicit sector \(\perp\) induced replacement;
 \item readiness gates, such as evolution requiring LF-to-QCD matching;
 \item scheme compatibility;
 \item exact quantum-number closure;
 \item process factorization status;
 \item required convergence axes.
\end{itemize}

\subsection{Dependency closure}

The compiler first computes a least dependency closure
\begin{equation}
 \operatorname{Cl}(\Aset)
 =
 \bigcap
 \left\{
 B\supseteq\Aset\mid
 B\text{ is closed under all dependency rules}
 \right\}.
\label{eq:dependency-closure}
\end{equation}
Closure may add required implementation objects, not new physical assumptions.
For example, requesting a rank-two GTMD projection may add its projector,
Bessel convention, and OAM blocks to the plan.  It may not silently add a
missing physical Fock sector.

\subsection{Compilation states}

Compilation returns one of four states:
\begin{equation}
 \mathrm{Compile}(\Aset)=
 \begin{cases}
 \mathrm{REJECTED}(\mathfrak f),\\
 \mathrm{PARTIAL}(\mathfrak g),\\
 \mathrm{EXECUTABLE}(\Pplan),\\
 \mathrm{VALIDATED}(\Pplan,\mathfrak c),
 \end{cases}
\label{eq:compile-states}
\end{equation}
where \(\mathfrak f\) is a structured failure certificate,
\(\mathfrak g\) lists unresolved gates, and \(\mathfrak c\) is a validation
certificate.  A partial plan may support diagnostic calculations but cannot be
promoted to a physical prediction.

\begin{theorem}[Compilation soundness]
Assume that every plan node has a typed domain and codomain, every dependency
is present, every exclusion and replacement relation is satisfied, every
required readiness gate passes, and every operator/state identity is complete.
Then any execution path generated by the compiler is type-correct and contains
no pair of contributions declared equivalent-count-once or mutually exclusive
by the active provenance complex.
\end{theorem}

\begin{proof}
The compiler constructs the plan by topological ordering of the closed
dependency graph.  Domain/codomain compatibility is checked at each edge.
Before a sum or composition node is emitted, the active subgraph is reduced by
all replacement, exclusion, and two-cell count rules.  Therefore no forbidden
pair survives into the execution graph.  Induction over the ordered plan nodes
establishes type correctness.  The theorem is a statement about the compiler;
it does not prove that the physical Hamiltonian or factorization theorem is
accurate.
\end{proof}

\subsection{Failure certificates}

A failure certificate records:
\begin{equation}
 \mathfrak f=
 \left(
 \mathrm{rule\ id},
 \mathrm{objects},
 \mathrm{requested\ target},
 \mathrm{minimal\ conflicting\ set},
 \mathrm{repair\ classes},
 \mathrm{provenance\ path}
 \right).
\label{eq:failure-certificate}
\end{equation}
The repair classes are restricted to scientifically meaningful actions such
as adding a missing operator, selecting a different branch, completing a
matching calculation, or downgrading the requested status.  ``Fit a new TMD
coefficient'' is not a generic repair.

\section{Prediction plans and execution semantics}
\label{sec:prediction-plan}

\subsection{Typed execution DAG}

A \PredictionPlan\ is a directed acyclic graph
\begin{equation}
 \Pplan=(V_{\Pplan},E_{\Pplan},\tau,\pi),
\label{eq:plan-dag}
\end{equation}
where \(\tau\) assigns a typed computational object to every node and \(\pi\)
assigns provenance, readiness, approximation, and caching metadata.  Typical
nodes are:
\begin{equation}
\begin{aligned}
 \text{basis}&\to\text{Hamiltonian}\to\text{renormalization}\to
 \text{solver}\to\text{state}\\
 &\to\text{operator contraction}\to\text{GTMD}\to
 \text{TMD/process result}.
\end{aligned}
\label{eq:plan-chain}
\end{equation}
The graph may branch to compute several observables from one state.  Shared
nodes are evaluated once per microscopic member and reused by immutable
content identity.

\subsection{Execution semantics}

The semantic map
\begin{equation}
 \llbracket\cdot\rrbracket:
 \PlanCat\longrightarrow\mathbf{Result}
\label{eq:execution-semantics}
\end{equation}
assigns to every executable plan a result envelope containing numerical arrays,
uncertainties, residuals, source identities, assumption identity, and
validation status.  The map is deterministic for fixed random seeds, numerical
libraries, precision, and hardware policy, all of which belong to the
execution manifest.

\subsection{Readiness lattice}

Readiness is a product of independent statuses rather than one boolean flag:
\begin{equation}
 \mathfrak r=
 \left(
 r_{\mathrm{state}},
 r_{\mathrm{current}},
 r_{\GTMD},
 r_{\mathrm{Wilson}},
 r_{\mathrm{nuclear}},
 r_{\mathrm{match}},
 r_{\mathrm{evol}},
 r_{\mathrm{process}},
 r_{\mathrm{infer}}
 \right).
\label{eq:readiness-product}
\end{equation}
A plan may be state-ready but not GTMD-ready, or GTMD-ready but not
process-ready.  The compiler cannot infer a later status from an earlier one.

\subsection{Caching and reproducibility}

A cache key is
\begin{equation}
 \mathfrak h_v=
 H\left(
 \mathrm{node\ type},
 \mathrm{inputs},
 \mathrm{source\ hashes},
 \Aset,
 \mathrm{numerical\ manifest}
 \right).
\label{eq:cache-key}
\end{equation}
Cache reuse is allowed only when the full key matches.  Two tensors with equal
values but different operator, scheme, or provenance identity are not
interchangeable.

\section{Symmetry-adapted tensor spaces}
\label{sec:symmetric-tensors}

\subsection{Representation and multiplicity decomposition}

Let \(G\) denote a compact exact internal or kinematic symmetry represented on
an edge space \(V_e\).  At finite truncation,
\begin{equation}
 V_e
 \cong
 \bigoplus_{q\in\widehat G_e}
 \left(
 D_{e,q}\otimes V_q
 \right),
\label{eq:edge-rep-decomposition}
\end{equation}
where \(V_q\) is an irreducible representation space and \(D_{e,q}\) is its
degeneracy or multiplicity space.  Quantum numbers in \(q\) may combine exact
charges from several commuting groups.  Representation dimensions are not
bond dimensions: the adjustable variational freedom resides in the degeneracy
spaces.

For an invariant tensor
\(T:\bigotimes_{i=1}^{n}V_{e_i}\rightarrow V_{e_0}\), equivariance requires
\begin{equation}
 T\left[\bigotimes_i U_{e_i}(g)\right]
 =U_{e_0}(g)T,
 \qquad g\in G.
\label{eq:equivariance}
\end{equation}
The tensor decomposes into structural intertwiners and reduced tensors,
\begin{equation}
 T=
 \bigoplus_{\bm q,\alpha}
 P_{\bm q,\alpha}
 \otimes
 Q_{\bm q,\alpha},
\label{eq:symmetric-tensor-decomposition}
\end{equation}
where \(Q_{\bm q,\alpha}\) contains Clebsch--Gordan and recoupling data fixed by
symmetry, while \(P_{\bm q,\alpha}\) contains the dynamical variational
parameters.  This is the tensor-network form of the Wigner--Eckart principle
\cite{Singh2009,Weichselbaum2012}.

\begin{proposition}[Symmetry closure under contraction]
The contraction of equivariant tensors over a common representation-typed edge
is equivariant.  Consequently, a network assembled only from admissible
intertwiners lies in the requested exact symmetry sector without a numerical
penalty term.
\end{proposition}

\begin{proof}
The group action on the contracted edge appears once in a representation and
once in its dual.  Their contraction is invariant.  Applying
\cref{eq:equivariance} to the remaining external edges yields the equivariance
condition for the composite tensor.
\end{proof}

\subsection{Multiple non-Abelian symmetries}

The implementation may use a product of representation labels such as
\begin{equation}
 q=
 \left(
 R_c,
 I,
 S,
 \Jz,
 B,
 Q,
 \mathbb Z_2^{F},
 \ldots
 \right),
\label{eq:compound-charge}
\end{equation}
provided unlike structures are not falsely promoted to one exact global group.
Color is a gauge constraint, isospin may be approximate, and light-front
rotations have kinematic and dynamical components.  The compound charge is a
software label for compatible block structure, not a claim that all entries
have the same physical status.

For non-Abelian factors with outer multiplicity, the label \(\alpha\) in
\cref{eq:symmetric-tensor-decomposition} is mandatory.  C7/H0 already shows why:
there are two independent \(qqqg\) singlets and three independent
\(qqqq\bar q\) singlets.  Dropping multiplicity labels erases physical color
channels even if the total singlet label is retained.

\subsection{Fermionic grading}

Fermionic statistics are represented by a \(\mathbb Z_2\)-graded tensor
category or an equivalent exact signed permutation implementation.  Swapping
homogeneous objects \(x,y\) produces
\begin{equation}
 \sigma(x\otimes y)
 =(-1)^{|x||y|}y\otimes x.
\label{eq:graded-swap}
\end{equation}
Every network reordering, fusion-tree transformation, and contraction-path
optimization must preserve this sign.  A bosonic contraction optimizer may be
used only after all fermionic swaps have been made explicit in the graph.

\subsection{Exact, approximate, and broken symmetries}

The compiler classifies each symmetry as:
\begin{equation}
 \mathrm{status}(G)=
 \mathrm{EXACT}\,|\,
 \mathrm{CONTROLLED\_BREAKING}\,|\,
 \mathrm{DIAGNOSTIC}\,|\,
 \mathrm{UNAVAILABLE}.
\label{eq:symmetry-status}
\end{equation}
Exact symmetries define absent blocks.  Controlled breaking adds explicit
operator blocks with a declared power counting.  Diagnostic symmetries supply
residuals but do not prune the basis.  This distinction is essential for
isospin breaking and rotational restoration.

\section{Physical tensor-edge types}
\label{sec:edge-types}

\subsection{No anonymous indices}

A state-network edge is typed by
\begin{equation}
 e=
 \left(
 \begin{gathered}
 \nu,
 a,
 R_c,
 R_f,
 \lambda,
 \Lz,
 \Jz,
 \ell_T,
 \mathbb Z_2^F,\\
 x\text{-mode},
 n,m\text{ transverse modes},
 \alpha,
 r,
 \mathfrak s,
 \mathrm{orientation}
 \end{gathered}
 \right).
\label{eq:full-edge-type}
\end{equation}
Here \(\nu\) is the Fock or nuclear sector, \(a\) is the parton species,
\(R_c,R_f\) are color and flavor representations, \(\lambda\) is parton
helicity, \(\ell_T\) is a transverse tensor or harmonic label, \(\alpha\) is a
basis or representation multiplicity, \(r\) is the regulator/truncation
identity, and \(\mathfrak s\) carries any operator scheme data.  The orientation
indicates whether the edge is incoming, outgoing, dual, bra, ket, or virtual.

\begin{requirement}[Physical index ledger]
Every tensor index must map to a physical degree of freedom, an exact
representation multiplicity, or a declared numerical compression index.  An
untyped array axis is not admissible in the central network.
\end{requirement}

\subsection{Direct sums over Fock sectors}

The microscopic state is a direct sum, not one dense tensor with a sector
label appended:
\begin{equation}
 |\Psi_N\rangle
 =
 \bigoplus_{\nu\in\cF_r}
 |\psi_{N,\nu}\rangle,
 \qquad
 \cF_r=\{qqq,qqqg,qqqq\bar q,\ldots\}.
\label{eq:fock-direct-sum}
\end{equation}
Each sector may have its own internal tree topology and symmetry blocks.
Sector-changing Hamiltonian or operator terms are explicit maps between the
corresponding network objects.  An absent sector is structurally zero; it is
not represented by a fitted distribution attached downstream.

\subsection{Momentum constraints}

A Fock-sector leaf set carries
\begin{equation}
 x_i>0,
 \qquad
 \sum_i x_i=1,
 \qquad
 \sum_i\bm\kappa_{iT}=0.
\label{eq:momentum-constraints-tn}
\end{equation}
The tensor network may use a constrained basis, Jacobi coordinates, or an
explicit projector, but the constraint method is part of its identity.  A
virtual bond cannot absorb a momentum mismatch without a declared recoil or
matching map.

\subsection{State edges versus operator edges}

Operator networks carry additional labels:
\begin{equation}
 e_{\cO}=
 \left(
 e_{\mathrm{state}},
 \Gamma,
 P_{\mathrm{in}},P_{\mathrm{out}},
 \gamma,
 R_{\mathrm{link}},
 \mu,\zeta,
 \text{map class}
 \right).
\label{eq:operator-edge-type}
\end{equation}
A state tensor cannot be connected directly to an operator edge unless the
momentum fiber, representation, regulator, Wilson-path, and scheme identities
match or an explicit adapter is present.

\section{Gauge, color, and Wilson-path structure}
\label{sec:gauge-color}

\subsection{Physical singlets and virtual representations}

The hadron state satisfies
\begin{equation}
 Q^a_{\mathrm{color}}|\Psi_h\rangle=0.
\label{eq:color-singlet-tn}
\end{equation}
This is enforced by color intertwiners at the network nodes.  Virtual edges may
carry nontrivial color representations and multiplicities, but every open
physical hadron boundary is projected to the declared singlet.  Gauge-invariant
tensor-network constructions in lattice gauge theory provide a useful
computational analogy: local constraints should be built into the tensor ansatz
rather than restored by a penalty after optimization
\cite{Tagliacozzo2014,Canals2024}.

\subsection{Network gauge is not color gauge}

For any internal bond, an invertible matrix \(X_e\) may be inserted with its
inverse on the neighboring tensor without changing the represented state:
\begin{equation}
 A\longmapsto AX_e,
 \qquad
 B\longmapsto X_e^{-1}B.
\label{eq:virtual-gauge}
\end{equation}
This is tensor-network parameter redundancy.  It is unrelated to the local
QCD gauge transformation of quark, gluon, and Wilson-line fields.  Canonical
forms fix \cref{eq:virtual-gauge}; they do not fix light-front gauge or determine
a Wilson line.  The geometric structure of tensor-network gauge orbits can be
used to define stable tangent spaces and optimization metrics
\cite{Haegeman2014}.

\subsection{Wilson paths remain operator data}

An ordered gluon link pair
\begin{equation}
 \left(\gamma_1,\gamma_2,R_A,\cC_f\text{ or }\cC_d\right)
\label{eq:ordered-link-network}
\end{equation}
is represented by operator-network nodes and path-typed edges.  It is not
encoded in the topology of the state TTN.  Future/past reversal acts on the
operator identity, endpoints, complex conjugation, and helicity phases; it is
not a reversal of the TTN drawing.

\subsection{Gauge-compatible contraction theorem}

\begin{proposition}[Color closure]
Let every state node and operator node be an admissible color intertwiner and
let every contracted color edge pair a representation with its dual.  If all
external hadron edges are singlets, the resulting scalar or hadron matrix
element is color invariant.
\end{proposition}

\begin{proof}
The tensor product of intertwiners is an intertwiner, and contraction with the
dual evaluation map is invariant.  After all internal contractions, the only
external representation is the singlet, on which the color action is trivial.
\end{proof}

\section{Network topology as a declared approximation}
\label{sec:network-topology}

\subsection{Tree ansatz for the first microscopic state}

The first state ansatz is a symmetry-adapted tree tensor network
\begin{equation}
 \cT=(V_{\cT},E_{\cT}),
 \qquad
 \text{acyclic and connected},
\label{eq:tree-graph}
\end{equation}
whose leaves are physical basis factors and whose internal nodes are fusion or
variational tensors.  Tree networks admit exact canonical forms and efficient
contraction when entanglement across each edge is controlled
\cite{Shi2006,Gunst2018}.

A representative three-quark fusion tree is
\begin{equation}
 \left(q_1\otimes q_2\right)
 \xrightarrow{\ C_{12}\ }
 \left(R_{12},I_{12},S_{12},L_{12},\alpha_{12}\right),
 \qquad
 \left(R_{12},\ldots\right)\otimes q_3
 \xrightarrow{\ C_{(12)3}\ }
 \left(\mathbf 1_c,I,\Jz\right).
\label{eq:qqq-fusion-tree}
\end{equation}
The structural couplers and the dynamical reduced tensors remain separate.

\subsection{Topology choices}

Several topologies may represent the same finite Hilbert space:
\begin{itemize}
 \item parton-first trees, fusing particle slots before orbital modes;
 \item color-first or spin-first fusion trees;
 \item momentum-cluster trees, grouping strongly coupled basis regions;
 \item sector trees, joining \(qqq\), \(qqqg\), and sea blocks through explicit
 sector-changing maps;
 \item operator-adapted trees optimized for a selected family of matrix
 elements.
\end{itemize}
At full bond support these are changes of representation.  At truncated bond
dimension they define different variational manifolds and are therefore
assumption branches or diagnostic alternatives.

\begin{definition}[Topology identity]
A topology identity contains the graph, leaf ordering, fusion order,
representation convention, edge orientations, fermion-swap record, and
canonical-gauge convention.  Two diagrams that look similar but differ in any
of these fields are not the same topology identity.
\end{definition}

\subsection{Hyperedges and interaction networks}

Hamiltonian vertices and multi-field operators are naturally represented as
hyperedges or high-valence tensors.  They need not be forced into the state
fusion tree.  A structured open-system viewpoint is useful: state and operator
subnetworks expose typed interfaces and compose only along matching boundaries,
analogous to structured or decorated cospans
\cite{Fong2015,BaezCourser2020,Courser2020}.

\subsection{Loops and higher-dimensional ansatz classes}

A looped tensor network or \PEPS-like ansatz may eventually be useful for
sector-rich states, but it introduces more difficult gauge fixing,
contraction, and canonical-form problems.  It is not adopted merely because it
is more expressive.  The compiler may expose
\begin{equation}
 \mathrm{topology\ class}
 =\TTN\,|\,\MPS\,|\,\PEPS\,|\,\mathrm{general\ hypergraph},
\label{eq:topology-class}
\end{equation}
with independent readiness and contraction requirements.

\section{Exact tensorization and canonical gauges}
\label{sec:exact-tensorization}

\subsection{Schmidt support on a tree}

Cutting an edge \(e\) of a tree partitions the leaves into \(L_e\) and \(R_e\).
Every finite state admits a Schmidt decomposition
\begin{equation}
 |\Psi\rangle
 =
 \sum_{s=1}^{r_e}
 \sigma_{e,s}
 |L_{e,s}\rangle|R_{e,s}\rangle,
\label{eq:schmidt-edge}
\end{equation}
resolved further into symmetry sectors:
\begin{equation}
 \rho_{L_e}
 =
 \bigoplus_q
 \rho_{L_e}^{(q)}\otimes\Id_{V_q}.
\label{eq:symmetry-schmidt}
\end{equation}
The virtual edge must contain every \((q,s)\) with nonzero support for an exact
representation.

\begin{theorem}[Exact TTN representation]
Let \(|\Psi\rangle\) be a finite-dimensional state and \(\cT\) a fixed tree on
its physical factors.  If the bond space on every edge \(e\) has dimension at
least the Schmidt rank \(r_e\), with all required symmetry sectors and
multiplicities retained, then there exists a TTN on \(\cT\) representing
\(|\Psi\rangle\) exactly.
\end{theorem}

\begin{proof}
Choose a root and recursively perform Schmidt decompositions from the leaves
toward the root.  Each decomposition defines an isometry from the child spaces
to the retained Schmidt space.  The assumed bond dimensions contain the full
support, so no singular value is discarded.  Symmetry-resolved decompositions
preserve the structural intertwiners and factor only the degeneracy tensors.
Contraction of the resulting isometries reconstructs the original coefficient
tensor.
\end{proof}

\subsection{Canonical gauge}

For an oriented tree with orthogonality center \(c\), tensors away from \(c\)
are chosen isometric:
\begin{equation}
 \sum_{\mathrm{children},\mathrm{phys}}
 A_v^{\dagger}A_v=\Id_{e(v\to c)}.
\label{eq:tree-isometry}
\end{equation}
This makes the norm local to the center and stabilizes variational updates.
Canonicalization records the virtual gauge transformations applied on every
edge.  Degenerate Schmidt values require a deterministic basis convention or
a subspace-valued identity rather than arbitrary vectors.

\subsection{Phase and member identity}

The exact state and its tensorized representation share one phase convention.
Across resolutions or assumption branches, state identity is tracked by
projected overlaps, principal angles, current fingerprints, and symmetry
content.  Tensor canonicalization may not redefine the microscopic member
identity independently at each observable.

\subsection{Exact reconstruction gates}

An exact tensorization must reproduce:
\begin{equation}
 \begin{aligned}
 \norm{|\Psi_{\TN}\rangle-|\Psi_{\mathrm{exact}}\rangle}&=0,\\
 \langle\Psi_{\TN}|\cM^2|\Psi_{\TN}\rangle
 &=\langle\Psi_{\mathrm{exact}}|\cM^2|\Psi_{\mathrm{exact}}\rangle,\\
 \langle\Psi_{\TN}|\cO_i|\Psi_{\TN}\rangle
 &=\langle\Psi_{\mathrm{exact}}|\cO_i|\Psi_{\mathrm{exact}}\rangle
 \end{aligned}
\label{eq:exact-tn-gates}
\end{equation}
for a benchmark operator family containing charge, OAM-sensitive, color, and
current observables, not energy alone.

\section{Variational tree-tensor-network solver}
\label{sec:variational-ttn}

\subsection{Rayleigh--Ritz objective}

For a normalized tensor-network state, the variational energy is
\begin{equation}
 E[\Nstate]
 =
 \frac{
 \langle\Psi(\Nstate)|\cM^2|\Psi(\Nstate)\rangle
 }{
 \langle\Psi(\Nstate)|\Psi(\Nstate)\rangle
 }.
\label{eq:rayleigh-quotient}
\end{equation}
Optimization is restricted to the symmetry blocks and topology selected by the
assumption bundle.  A compressed exact state is an initialization or
representation benchmark; it is not a variational solver until
\cref{eq:rayleigh-quotient} is minimized using tensor updates.

\subsection{Local effective problem}

With all tensors except \(A_v\) fixed, the objective becomes
\begin{equation}
 E(A_v)
 =
 \frac{A_v^{\dagger}H_{\eff}^{(v)}A_v}
 {A_v^{\dagger}N_{\eff}^{(v)}A_v}.
\label{eq:local-generalized-eigen}
\end{equation}
In a canonical tree gauge, \(N_{\eff}^{(v)}=\Id\) for a one-center update.
The local optimum is the lowest admissible eigenvector of
\(H_{\eff}^{(v)}\), resolved by symmetry block.  Sweeps move the orthogonality
center through the tree and update every active tensor until energy, state,
and observable residuals stabilize.

\subsection{Variational bounds}

\begin{proposition}[Upper bound]
For a self-adjoint finite Hamiltonian \(\cM^2\), every normalized TTN state in
the declared physical subspace satisfies
\begin{equation}
 E_{\TTN}\geq E_0,
\label{eq:variational-upper-bound}
\end{equation}
where \(E_0\) is the exact ground-state eigenvalue in that finite subspace.
\end{proposition}

\begin{proposition}[Nested bond monotonicity]
If the TTN variational manifolds satisfy
\(\cV_{\chi}\subseteq\cV_{\chi'}\) for
\(\chi\preceq\chi'\), then the globally minimized energies obey
\begin{equation}
 E_{\chi'}\leq E_{\chi}.
\label{eq:bond-monotonicity}
\end{equation}
A numerical violation signals incomplete optimization, inconsistent symmetry
spaces, or nonnested bond policies.
\end{proposition}

\subsection{Optimization routes}

The implementation may combine:
\begin{enumerate}
 \item one-site or two-site tree sweeps;
 \item direct local eigensolves;
 \item Krylov updates using matrix-free environments;
 \item gradient or natural-gradient optimization;
 \item subspace expansion and controlled bond growth;
 \item exact-state initialization on small towers.
\end{enumerate}
Every route records initialization, sweep ordering, stopping rule, local and
global residuals, and whether the reported state is the best result over
multiple starts.

\subsection{Excited states and near degeneracy}

Excited or symmetry-partner states may be obtained by block targeting,
projected Hamiltonians, or penalty projectors:
\begin{equation}
 \cM_k^2
 =
 \cM^2+
 \sum_{j<k}\mu_j|\Psi_j\rangle\langle\Psi_j|.
\label{eq:excited-penalty}
\end{equation}
Near degeneracies are tracked as subspaces.  The compiler must not assign a
persistent physical label solely from eigenvalue order.

\subsection{Tensor-network solver certificate}

A solver certificate contains:
\begin{equation}
 \mathfrak c_{\TN}=
 \left(
 \begin{gathered}
 \text{topology},\ \text{symmetry blocks},\ \chi_e,\
 \text{canonical gauge},\\
 E,\ \norm{\cM^2\Psi-E\Psi},\
 \text{sweep history},\ \text{discarded weights},\\
 \text{exact/Krylov comparison},\
 \text{observable residuals},\
 \text{source and hardware hashes}
 \end{gathered}
 \right).
\label{eq:tn-solver-certificate}
\end{equation}
Energy convergence without current, OAM, or state-fidelity convergence is
insufficient.

\section{Bond spaces, entanglement, and adaptive compression}
\label{sec:bond-adaptation}

\subsection{Symmetry-resolved Schmidt spectra}

For each edge \(e\), let
\begin{equation}
 \{\sigma_{e,q,s}\}
\label{eq:symmetry-schmidt-values}
\end{equation}
be the singular values in charge sector \(q\).  A truncation retains
\(\chi_{e,q}\) values in each sector and has discarded weight
\begin{equation}
 \eps_e
 =
 \sum_q\sum_{s>\chi_{e,q}}
 \sigma_{e,q,s}^{2}.
\label{eq:discarded-weight}
\end{equation}
The total bond dimension is
\begin{equation}
 \chi_e=\sum_q \chi_{e,q}\dim V_q,
\label{eq:physical-bond-dimension}
\end{equation}
but computational storage is controlled primarily by the degeneracy
dimensions \(\chi_{e,q}\).

\subsection{Required-sector floor}

A bond policy first allocates a required floor
\begin{equation}
 \chi_{e,q}\geq \chi_{e,q}^{\min}
\label{eq:required-sector-floor}
\end{equation}
for representation sectors needed by exact constraints or requested
observables.  Remaining resources may be allocated by Schmidt weight,
observable sensitivity, or error estimates.  A small singular value does not
justify deleting the only block capable of supporting a requested OAM or color
interference without first downgrading the prediction plan.

\subsection{Adaptive objective}

For a resource budget \(B\), one possible allocation solves
\begin{equation}
 \min_{\{\chi_{e,q}\}}
 \left[
 \sum_e w_e\eps_e
 +\sum_i \lambda_i\delta\cO_i^2
 \right]
 \quad\text{subject to}\quad
 \operatorname{cost}(\{\chi_{e,q}\})\leq B.
\label{eq:adaptive-bond-objective}
\end{equation}
The observable terms prevent energy-dominated compression from erasing
spin--orbit, tensor, or GTMD-sensitive directions.  The weights and target
observables belong to the assumption bundle.

\subsection{Entanglement diagnostics}

The edge entropy is
\begin{equation}
 S_e=-\sum_{q,s}\sigma_{e,q,s}^{2}
 \ln\sigma_{e,q,s}^{2}.
\label{eq:edge-entropy}
\end{equation}
It is a diagnostic of the chosen bipartition, not a universal physical
observable.  Comparison across topologies requires a map between cuts.  High
entropy or slow Schmidt decay identifies where tree topology or bond support
may be inadequate.

\subsection{Independent network truncation uncertainty}

Tensor-network uncertainty is reported separately from basis and Fock
uncertainty:
\begin{equation}
 \delta_{\TN}\cO
 =
 \cO(\chi,\cT)-
 \cO(\chi',\cT')
\label{eq:tn-truncation-error}
\end{equation}
for nested bond dimensions and controlled topology variations.  It may not be
absorbed into Hamiltonian EFT discrepancy or posterior parameter uncertainty.

\begin{principle}[Observable convergence]
Bond convergence is claimed for a set of matched observables and state
residuals, not from the energy or discarded weight alone.
\end{principle}

\section{Tensor operators and contraction planning}
\label{sec:operator-networks}

\subsection{Hamiltonian operator network}

The Hamiltonian is a sum of typed operator subnetworks,
\begin{equation}
 \cM^2=
 \sum_{a=1}^{N_{\mathrm{term}}}
 c_a\cO_a,
\label{eq:hamiltonian-sum-network}
\end{equation}
where each \(\cO_a\) retains source/target sector, symmetry, regulator,
Hermitian partner, and parameter ownership.  The operator representation may
be a tree tensor operator, a sum of factored terms, or a general hypergraph.
The state topology does not force the Hamiltonian into the same topology.

\subsection{Shared environments}

For several observables evaluated on one state, contraction environments are
reused only when their tensor and operator identities match.  A common
left/right or tree environment reduces repeated cost while preserving one
microscopic member.  Shared environments are especially useful for the family
of helicity and transverse-rank projectors derived from one GTMD parent.

\subsection{Off-forward matrix elements}

A GTMD contraction joins different momentum fibers:
\begin{equation}
 W_{\Lambda'\Lambda}^{[\Gamma,\gamma]}
 (x,\kT,\DT)
 =
 \cC\left[
 {\Nstate}_{P',\Lambda'}^{\dagger},
 {\Nop}_{\Gamma,\gamma}(P',P),
 {\Nstate}_{P,\Lambda}
 \right].
\label{eq:gtmd-tn-contraction}
\end{equation}
The bra and ket may share a state bundle but use the exact symmetric recoil map
from Volume~II.  Their canonical gauges need not be identical; the contraction
plan records the transition and phase conventions.

\subsection{Density and completely positive contractions}

Forward Gram correlators may be evaluated by contracting an amplitude network
with its conjugate over unresolved states.  A completely positive nuclear or
detector map is represented only after a justified partial trace.  The
double-layer network is therefore distinct from a Wigner quasidistribution or
an off-forward GTMD, for which pointwise positivity is not required.

\subsection{Contraction-order optimization}

For irregular operator networks, contraction order can alter computational
cost by orders of magnitude while leaving the scalar invariant.  The
\ContractionPlan\ stores a hypergraph path, slicing policy, memory estimate,
fermionic swap ledger, and deterministic optimizer seed.  Hypergraph-based
contraction-path optimization is an admissible numerical tool
\cite{GrayKourtis2021}; it is not a physical model choice unless an
approximate contraction is introduced.

\subsection{Approximate contraction}

If contraction is truncated, sampled, or distributed with finite tolerance,
its approximation record contains:
\begin{equation}
 \mathfrak a_{\cC}=
 \left(
 \text{method},\text{error estimator},\text{bias test},
 \text{resource budget},\text{seed},\text{validation oracle}
 \right).
\label{eq:approx-contraction-record}
\end{equation}
Approximate contraction uncertainty is numerical, separate from the state
ansatz and physical truncation.

\section{Common-parent GTMD and TMD predictions}
\label{sec:tmd-predictions}

\subsection{The tensor-network parent chain}

For one microscopic member and assumption bundle,
\begin{equation}
 \Aset
 \longmapsto
 |\Psi_N(\Aset)\rangle
 \longmapsto
 W_{q/N},W_{\bar q/N},W_{g/N}
 \longmapsto
 \{F_A\}
\label{eq:tn-parent-chain}
\end{equation}
is evaluated by one state network and a family of typed operator networks.
The named functions are coordinates on the parent correlator:
\begin{equation}
 F_A^{a/N}
 =
 \Tr\left[
 P_A^{a}
 W_{a/N}
 \right].
\label{eq:named-tmd-projection}
\end{equation}
No tensor-network parameter is owned by a named TMD unless the corresponding
operator is explicitly present in the Hamiltonian or current basis.

\subsection{Common reductions}

The same off-forward contraction must support
\begin{equation}
 \begin{aligned}
 \mathsf T[W]&=W\big|_{\DT=0},\\
 \mathsf G[W]&=\int\dd^2\kT\,W,\\
 \mathsf F[W]&=\int\dd x\,\dd^2\kT\,W,\\
 \mathsf W[W]&=\int\frac{\dd^2\DT}{(2\pi)^2}
 e^{-i\bD\cdot\DT}W.
 \end{aligned}
\label{eq:common-reductions-tn}
\end{equation}
These are regulated reductions until the matching objects of Volume~V are
applied.  Closure failure is traced to the state, operator, recoil, matching,
or numerical network rather than repaired independently in each output.

\subsection{Transverse rank}

A projection request carries its harmonic \(m\), Bessel order \(J_{|m|}\),
extracted mass powers, and Fourier phase.  The compiler uses these labels to
select the required OAM and helicity blocks before contraction.  If the current
state network lacks a required block, the result is unavailable rather than
zero by default.

\subsection{T-even and link-odd routes}

At zeroth Wilson order, a real Hamiltonian state supplies the T-even common
parent.  Link-odd projections require the Volume~III rescattering operator
network and its cut ledger:
\begin{equation}
 W_{\mathrm{odd}}
 =
 \frac12
 \left(
 W^{[+]}-\Theta^{-1}W^{[-]}\Theta
 \right).
\label{eq:link-odd-tn}
\end{equation}
The same state tensors are reused, but the operator and path networks differ.
A universal scalar phase multiplying all spin projectors is not an admissible
network node.

\subsection{Conditional prediction object}

A prediction is stored as
\begin{equation}
 \mathfrak F_A=
 \left(
 F_A,
 \Aset,
 \text{member},
 \text{operator id},
 \text{network ids},
 \text{residuals},
 \text{matching status},
 \text{readiness}
 \right).
\label{eq:conditional-prediction-object}
\end{equation}
The phrase ``the model predicts \(F_A\)'' is incomplete unless the assumption
bundle and readiness status are supplied.

\section{Open interfaces to Wilson, nuclear, matching, and process layers}
\label{sec:open-interfaces}

\subsection{Why one category is insufficient}

The state tensor network lives in the amplitude layer \(\Amp\); density and
partial-trace maps live in \(\Dens\); regulator and scheme conversions live in
\(\Match\); GTMD/TMD/GPD/current reductions live in \(\Red\); and hard-process
assembly lives in \(\Proc\).  The prediction compiler coordinates these layers
without identifying their morphisms.

An open subsystem exposes a typed boundary.  Schematically,
\begin{equation}
 X_{\mathrm{in}}
 \longrightarrow
 N
 \longleftarrow
 X_{\mathrm{out}},
\label{eq:structured-cospan}
\end{equation}
where the apex \(N\) carries the internal network and its assumption
``decoration.''  Structured and decorated cospan formalisms motivate this
interface discipline, while double categories distinguish system composition
from transformations between systems
\cite{Fong2015,Courser2017,BaezCourser2020,Courser2020}.

\subsection{Wilson interface}

The Wilson plugin consumes:
\begin{equation}
 \left(
 \text{microscopic state bundle},
 \text{gauge-field vertex basis},
 \text{path identity},
 \text{cut and pole policy},
 \text{Wilson order}
 \right)
\label{eq:wilson-plugin-input}
\end{equation}
and returns an operator network with a phase budget.  It does not modify the
state topology silently.  Higher Wilson order may require new Fock sectors,
in which case the compiler must recompile the state plan rather than insert a
larger phase into the old state.

\subsection{Nuclear interface}

The nuclear plugin consumes correlated proton and neutron microscopic members,
complete helicity matrices, matched nucleon operators, and shared phase/soft
metadata.  It returns a spin-1 nuclear state network and nuclear operator
network.  Partonic Wilson staples and coherent nuclear rescattering retain
separate interfaces and an explicit overlap relation when they share a soft
region.

\subsection{Matching and evolution interface}

The matching plugin requires a closed regulated operator basis and produces
\begin{equation}
 \widetilde{\bm W}^{\QCD}
 =
 \bm Z_{\LF\to\QCD}
 \otimes_x
 \widetilde{\bm W}^{\LF}
 +\bm\Delta_{\trunc}.
\label{eq:matching-plugin}
\end{equation}
The result remains attached to the originating state-network member.  Evolution
acts on the matched operator coordinates, not on arbitrary internal TTN bonds.

\subsection{Process interface}

A process plugin supplies hard factors, fragmentation functions, link and
color maps, soft prescription, \(W+Y\) accuracy, and factorization/Glauber
status.  A TMD tensor network cannot be consumed by a process marked
\texttt{BROKEN} or \texttt{UNASSESSED} as though universal factorization held.

\section{Provenance two-complex and topological audit}
\label{sec:provenance-complex}

\subsection{Cells and physical meaning}

For an assumption bundle \(\Aset\), construct a finite cellular complex
\(\Kprov(\Aset)\):
\begin{itemize}
 \item \(0\)-cells are typed states, operators, schemes, networks, and results;
 \item oriented \(1\)-cells are derivations, transitions, reductions,
 matching maps, replacements, or execution steps;
 \item \(2\)-cells are attached only when a physical identity declares two
 paths equivalent, subtractive, replacing, or count-once.
\end{itemize}
The chain condition is
\begin{equation}
 C_2(\Kprov)
 \xrightarrow{\partial_2}
 C_1(\Kprov)
 \xrightarrow{\partial_1}
 C_0(\Kprov),
 \qquad
 \partial_1\partial_2=0.
\label{eq:prov-chain}
\end{equation}

\subsection{Examples of two-cells}

The required cells include:
\begin{enumerate}
 \item explicit \(qqqg\) propagation versus the induced operator obtained by
 eliminating that sector;
 \item a cut represented in an eikonal denominator versus the same physical
 on-shell support represented in a light-front resolvent;
 \item unsubtracted collinear overlap versus the half-soft subtraction of the
 same rapidity region;
 \item explicit nuclear \(NN\pi\) content versus pion dressing already
 contained in the matched nucleon state;
 \item alternative tensor topologies that are exactly gauge-equivalent at full
 bond support but distinct at finite rank.
\end{enumerate}
A two-cell may carry a coefficient, sign, matching remainder, tolerance, and
validity domain.

\subsection{Path quotient and count-once semantics}

Let \(\cP(s,t)\) denote formal source-to-target paths.  Declared two-cells
generate an equivalence/subtraction relation \(\sim\).  The composition engine
works on the quotient presentation
\begin{equation}
 \cP_{\mathrm{phys}}(s,t)
 =
 \cP(s,t)/{\sim}
\label{eq:path-quotient}
\end{equation}
together with additive and exclusive branch labels.

\begin{proposition}[Count-once reduction]
If two active paths represent the same contribution and a valid count-once
cell identifies them, then the compiled observable contains one representative
of the equivalence class plus any declared matching remainder, independent of
which representative is selected for numerical execution.
\end{proposition}

\begin{proof}
The cell boundary is \(p_1-p_2-r\), where \(r\) is an optional remainder path.
The compiler normalizes the active graph by the corresponding rewrite rule
before emitting an additive node.  Therefore either \(p_1\) or \(p_2+r\) is
present, but not both.  Equality holds within the cell's declared domain and
tolerance.
\end{proof}

\subsection{Homology as an obstruction signal}

The group
\begin{equation}
 H_1(\Kprov)
 =
 \Ker\partial_1/\operatorname{im}\partial_2
\label{eq:prov-homology-v8}
\end{equation}
identifies closed provenance cycles not resolved by existing two-cells.  A
nontrivial class may indicate a missing subtraction, unresolved alternative,
unitarity relation, or genuinely distinct interference mechanism.  It is an
audit result, not a topological value of the physical amplitude.

\subsection{Higher cells}

Three-cells may eventually express coherence among several two-cell relations,
such as compatibility of sector elimination, operator matching, and soft
subtraction.  They are introduced only when independent pairwise cells produce
a concrete coherence ambiguity.  Decorative higher category theory is not a
deliverable.

\section{Differentiable networks and microscopic sensitivities}
\label{sec:differentiation}

\subsection{Differentiation through contractions}

For network parameters \(\theta_a\), automatic differentiation gives
\begin{equation}
 J_{ia}
 =
 \frac{\partial\cO_i}{\partial\theta_a}
\label{eq:tn-jacobian}
\end{equation}
through tensor contractions, provided decomposition and gauge-fixing steps have
stable derivative rules.  Differentiable tensor-network programming can
provide efficient gradients and higher derivatives
\cite{Liao2019}.

\subsection{Eigenstate derivatives}

For a nondegenerate normalized eigenstate,
\begin{equation}
 \cM^2(\theta)|\Psi_n\rangle
 =E_n|\Psi_n\rangle,
\label{eq:eigenproblem-derivative}
\end{equation}
Hellmann--Feynman gives
\begin{equation}
 \frac{\partial E_n}{\partial\theta_a}
 =
 \langle\Psi_n|
 \frac{\partial\cM^2}{\partial\theta_a}
 |\Psi_n\rangle.
\label{eq:hellmann-feynman}
\end{equation}
The orthogonal state derivative satisfies
\begin{equation}
 Q_n
 \left(\cM^2-E_n\right)
 Q_n
 |\partial_a\Psi_n\rangle
 =
 -Q_n
 \left(
 \partial_a\cM^2-\partial_a E_n
 \right)
 |\Psi_n\rangle,
\label{eq:implicit-state-derivative}
\end{equation}
where \(Q_n=\Id-|\Psi_n\rangle\langle\Psi_n|\).  Near degeneracies require
subspace derivatives rather than \cref{eq:implicit-state-derivative} for an
arbitrary eigenvector.

\subsection{Gauge-aware tangent vectors}

Because tensor parameters contain virtual gauge redundancy, raw parameter
derivatives are not unique.  The implementation chooses a horizontal tangent
gauge or canonical-center convention so that
\begin{equation}
 \delta|\Psi\rangle
 =
 \sum_v
 \frac{\partial|\Psi\rangle}{\partial A_v}\delta A_v
\label{eq:tn-tangent}
\end{equation}
contains no pure virtual-gauge direction.  This is the practical role of the
tensor-manifold geometry discussed for matrix-product states
\cite{Haegeman2014,Haegeman2013}.

\subsection{Differentiating SVD and truncation}

A singular-value crossing or exact degeneracy makes individual singular
vectors nondifferentiable.  The compiler therefore records block-subspace
projectors and either:
\begin{itemize}
 \item differentiates the invariant retained subspace;
 \item freezes the truncation pattern within a local neighborhood;
 \item uses a smooth spectral filter for a diagnostic route;
 \item marks the derivative unavailable.
\end{itemize}
It may not return an unstable gradient without a degeneracy warning.

\subsection{Adjoint and covariance propagation}

For parameter covariance \(C_{\theta}\), the linearized observable covariance
is
\begin{equation}
 C_{\cO}^{\mathrm{param}}
 =JC_{\theta}J^{\mathsf T}.
\label{eq:jacobian-covariance}
\end{equation}
The same member identity and tensor network are used for all observables.  For
large parameter sets, adjoint contractions compute selected gradients without
forming the full Jacobian.

\section{Conditional predictions and uncertainty semantics}
\label{sec:conditional-uncertainty}

\subsection{Within-assumption prediction}

For a fixed assumption bundle,
\begin{equation}
 p(F_A\mid\Aset,\cD_{\calib})
 =
 \int\dd\Theta\,\dd\nu\,
 p(F_A\mid\Theta,\nu,\Aset)
 p(\Theta,\nu\mid\Aset,\cD_{\calib}),
\label{eq:conditional-posterior-prediction}
\end{equation}
where \(\Theta\) contains shared microscopic parameters and \(\nu\) contains
controlled truncation and numerical choices.  The posterior does not contain
an independent parameter for every named TMD.

\subsection{Between-assumption uncertainty}

If a declared model-comparison rule assigns weights to complete branches, the
law of total covariance gives
\begin{equation}
 \Cov(F\mid\cD)
 =
 \mathbb E_{\Aset}
 \left[
 \Cov(F\mid\Aset,\cD)
 \right]
 +
 \Cov_{\Aset}
 \left[
 \mathbb E(F\mid\Aset,\cD)
 \right].
\label{eq:total-covariance-assumptions}
\end{equation}
The second term is structural branch variation.  It is reported separately
unless a scientifically justified model-averaging procedure is declared.
Mutually exclusive mechanisms cannot be combined component by component merely
to widen a band.

\subsection{Tensor-network uncertainty axes}

At minimum, report:
\begin{enumerate}
 \item bond-dimension truncation;
 \item topology choice;
 \item variational optimization residual;
 \item contraction approximation;
 \item basis and Fock truncation;
 \item Hamiltonian EFT and regulator uncertainty;
 \item Wilson, matching, evolution, nuclear, and process uncertainty;
 \item posterior parameter and data uncertainty.
\end{enumerate}
A small bond error does not imply a small Fock-sector error, and an exact
contraction does not imply physical completeness.

\subsection{Prediction versus scenario}

An output is a prediction relative to \(\Aset\) only if its normalization and
shape were not directly calibrated as a named observable.  Across assumption
branches, the scientifically correct language may be ``conditional prediction''
or ``model scenario.''  The result envelope stores this status explicitly.

\section{Compiler algorithm and proof objects}
\label{sec:compiler-algorithm}

\subsection{Compilation pipeline}

A reference compiler executes the following stages:
\begin{enumerate}
 \item canonicalize and hash the assumption bundle;
 \item load exact conventions and repository capabilities;
 \item compute dependency closure;
 \item solve branch, exclusion, replacement, and readiness constraints;
 \item construct the provenance subcomplex and normalize count-once paths;
 \item instantiate the physical Hilbert and symmetry-block tensor spaces;
 \item select state topology and bond policy;
 \item instantiate Hamiltonian, current, GTMD, Wilson, nuclear, matching, and
 process operator networks requested by the target;
 \item choose solver and contraction plans;
 \item emit validation gates and holdout observables;
 \item create an immutable prediction-plan identity.
\end{enumerate}

\subsection{Pseudocode}

\begin{verbatim}
def compile(bundle, target, repository):
    A = canonicalize_and_hash(bundle)
    capability = repository.capability_snapshot()
    A_closed = dependency_closure(A, capability)
    certificate = solve_admissibility(A_closed, target)
    if certificate.failed:
        return Rejected(certificate)
    K = build_provenance_complex(A_closed, target)
    K_normal = normalize_count_once_paths(K)
    state_spec = instantiate_state_network_spec(A_closed)
    operator_specs = instantiate_required_operators(A_closed, target)
    readiness = evaluate_readiness(state_spec, operator_specs, capability)
    plan = build_execution_dag(
        state_spec, operator_specs, K_normal, readiness, target
    )
    return Executable(plan) if readiness.sufficient else Partial(plan)
\end{verbatim}

\subsection{Compilation certificate}

A certificate records every rule evaluated and the minimal evidence used:
\begin{equation}
 \mathfrak c_{\mathrm{compile}}
 =
 \left(
 \Aset,
 \mathrm{target},
 \{R_i\mapsto s_i\},
 \Kprov,
 \mathrm{readiness},
 \Pplan,
 \mathrm{source\ hashes}
 \right).
\label{eq:compile-certificate}
\end{equation}
This permits independent replay without rerunning the physical solver.

\subsection{Semantic versioning}

Changing a numerical tolerance may create a new execution version without
changing the scientific assumption identity.  Changing a Fock sector,
Hamiltonian branch, Wilson order, tensor topology, matching scheme, or
calibration role creates a new scientific version.  The two version layers are
stored separately.

\section{Software architecture}
\label{sec:software}

\subsection{Core objects}

\begin{longtable}{L{0.27\textwidth}L{0.65\textwidth}}
\toprule
Object & Responsibility\\
\midrule
\endhead
\AssumptionBundle & Frozen scientific, regulator, algebraic, numerical, process, and inference choices with canonical identity.\\
\texttt{ChoiceAxis} & Singleton, exclusive, additive, nested, continuous, or diagnostic assumption semantics.\\
\texttt{AdmissibilityRule} & Typed dependency, exclusion, compatibility, readiness, or convergence predicate.\\
\texttt{CompilationCertificate} & Rule outcomes, minimal evidence, failure or readiness state, and plan identity.\\
\TensorEdgeType & Full physical/virtual representation, momentum, regulator, parity, and orientation identity.\\
\SymTensor & Structural intertwiners plus reduced degeneracy tensors; forbidden blocks absent.\\
\FusionTree & Leaf order, fusion order, recoupling convention, multiplicities, fermion swaps, and topology hash.\\
\texttt{BondSpace} & Per-edge, per-symmetry-sector degeneracy support and required-sector floors.\\
\texttt{BondDimensionPolicy} & Fixed, nested, adaptive, or observable-weighted allocation with cost and error records.\\
\StateTN & Canonical-gauge state tensors, microscopic member identity, topology, and solver certificate.\\
\OperatorTN & Hamiltonian, current, GTMD, Wilson, nuclear, matching, or process tensor network.\\
\texttt{VariationalTTNSolver} & Local effective solves, sweeps, subspace targeting, convergence, and multistart control.\\
\ContractionPlan & Hypergraph contraction order, slicing, memory, fermion swaps, approximation, and reproducibility.\\
\PredictionPlan & Typed execution DAG from assumptions and targets to result envelopes.\\
\ProvComplex & Typed graph, two-cells, path normalization, homology diagnostics, and count-once enforcement.\\
\texttt{ConditionalPrediction} & Result, assumption identity, member, operator, residual, covariance, and readiness.\\
\shortstack[l]{\texttt{CompilerCapability}\\\texttt{Snapshot}} & Repository objects, source hashes, supported schemes, solvers, and readiness gates.\\
\bottomrule
\end{longtable}

\subsection{Serialization}

All central objects require deterministic serialization, schema versioning,
source hashes, and round-trip tests.  Floating arrays are referenced by content
hash and immutable storage location rather than embedded redundantly in every
manifest.

\subsection{Distributed realization}

A distributed network stores ownership of tensor blocks, contraction tasks,
communication topology, deterministic reduction order, and precision.  Load
balancing may repartition blocks but cannot merge distinct physical identities.
Communication-aware contraction is a numerical optimization; its result must
match a serial oracle on smaller instances.

\subsection{Quantum simulation}

A quantum tensor-network or variational circuit backend is admissible only when
it represents the same typed Hilbert space, Hamiltonian, state topology, and
operator matrix elements and is benchmarked against exact or classical TTN
results.  A generic entangling circuit with no map to the microscopic tensor
blocks is outside this architecture.

\section{Formal requirements}
\label{sec:requirements}

The following identifiers are stable normative requirements for the first
implementation.  A requirement may be satisfied at validation-only scope
without implying physical prediction readiness.

\begin{longtable}{L{0.15\textwidth}L{0.77\textwidth}}
\toprule
ID & Requirement\\
\midrule
\endhead
V8.SEP.1 & State, operator/contraction, and provenance networks are distinct typed objects with different composition laws.\\
V8.SEP.2 & Tensor virtual gauge, QCD gauge, Wilson-path geometry, and provenance topology cannot be aliased.\\
V8.ASSUMP.1 & Every scientific run is generated from one frozen, content-addressed \AssumptionBundle.\\
V8.ASSUMP.2 & Choice axes distinguish required singleton, exclusive branch, additive mechanism, nested truncation, continuous parameter, and diagnostic variation.\\
V8.ASSUMP.3 & Refinement relations are declared only for nested choices; sibling physical branches remain incomparable.\\
V8.COMP.1 & Compilation computes dependency closure without silently adding missing physical dynamics.\\
V8.COMP.2 & Exclusion, replacement, scheme, readiness, factorization, and no-double-counting rules fail closed.\\
V8.COMP.3 & Rejected and partial compilations return structured minimal failure certificates.\\
V8.COMP.4 & Every executable plan has deterministic identity, source hashes, environment, and replayable rule outcomes.\\
V8.EDGE.1 & Every state and operator index has a complete \TensorEdgeType; anonymous axes are forbidden.\\
V8.SYM.1 & Exact symmetries are represented by invariant tensors and absent forbidden blocks, not penalties.\\
V8.SYM.2 & Structural intertwiners and reduced degeneracy tensors are stored separately.\\
V8.SYM.3 & Outer multiplicity labels are retained for every repeated irrep coupling.\\
V8.FERM.1 & Fermionic swaps and antisymmetry are exact and preserved by topology changes and contraction optimization.\\
V8.COLOR.1 & Every hadron state network is annihilated by the total color generator within tolerance.\\
V8.COLOR.2 & Ordered gluon links and \(f/d\) color contractions remain operator identities independent of state topology.\\
V8.TOP.1 & The state topology, fusion convention, leaf order, and canonical gauge form a versioned topology identity.\\
V8.TOP.2 & Finite-rank results from different topologies are distinct approximation branches.\\
V8.TN.1 & Exact tensorization reproduces the finite exact state and benchmark observables at full Schmidt support.\\
V8.TN.2 & A claimed variational solver minimizes the Rayleigh quotient and is not merely an SVD compression.\\
V8.TN.3 & Nested bond spaces give nonincreasing globally optimized variational energies.\\
V8.TN.4 & Bond manifests retain per-edge, per-symmetry-sector support, Schmidt spectra, discarded weight, and required-sector floors.\\
V8.TN.5 & Bond convergence is demonstrated for energy, residual, state identity, currents, OAM, and requested GTMD-sensitive observables.\\
V8.OP.1 & Hamiltonian, current, GTMD, Wilson, nuclear, matching, and process operator networks retain complete source/target and scheme identity.\\
V8.OP.2 & Off-forward contractions use typed incoming and outgoing momentum fibers and one authoritative recoil map.\\
V8.OP.3 & Shared contraction environments never merge different microscopic members or operator identities.\\
V8.PROV.1 & Every prediction has a provenance path from assumptions, state amplitudes, operators, matching, and numerical approximations.\\
V8.PROV.2 & Count-once, subtraction, replacement, and alternative relations are executable two-cells.\\
V8.PROV.3 & Nontrivial provenance homology is reported as an unresolved audit signal, not a physical invariant.\\
V8.GTMD.1 & Named TMDs are typed projections of common GTMD parents; named-function-specific latent coefficients are prohibited by default.\\
V8.GTMD.2 & TMD, GPD, PDF, current, and Wigner reductions share state, normalization, recoil, and operator ancestry.\\
V8.GTMD.3 & Missing OAM, helicity, Fock, or Wilson blocks produce an unavailable status rather than an assumed physical zero.\\
V8.OPEN.1 & Wilson, nuclear, matching/evolution, process, and inference layers compose through typed open interfaces and independent readiness gates.\\
V8.DIFF.1 & Automatic or adjoint derivatives carry tensor-gauge, eigenspace, decomposition, and truncation validity records.\\
V8.UNC.1 & Tensor-rank, topology, basis, Fock, physical, matching, evolution, nuclear, process, and statistical uncertainties remain separately identified.\\
V8.UNC.2 & Within-assumption posterior uncertainty is distinct from between-assumption structural variation.\\
V8.PRED.1 & Every result states whether it is a validation object, conditional prediction, fitted quantity, comparison, or physical process prediction.\\
V8.REGR.1 & New tensor-network and compiler packages cannot alter accepted production artifacts or registries unless a separate migration work package authorizes it.\\
\bottomrule
\end{longtable}

\section{Analytic and finite-dimensional benchmarks}
\label{sec:benchmarks}

\subsection{Benchmark TN-A: exact tensorization of the H0 basis}

Tensorize the C7 \(qqq\), \(qqqg\), and \(qqqq\bar q\) finite states on at
least two fusion trees.  At full bond support, reproduce the exact state,
color-generator residuals, antisymmetry, free mass, and reduced-vertex matrix
elements.  Recoupling between trees must be unitary within each complete
multiplicity space.

\subsection{Benchmark TN-B: structural versus degeneracy tensors}

Construct an invariant tensor with a repeated non-Abelian irrep channel.
Verify \cref{eq:symmetric-tensor-decomposition}, reconstruct the dense tensor,
and demonstrate that deleting the multiplicity label fails.  Repeat with a
fermionic swap crossing the fusion tree.

\subsection{Benchmark TN-C: canonical gauge}

Apply random invertible transformations on every virtual edge of an exact TTN.
Canonicalize to a deterministic center gauge and show state equality, norm
closure, identical observables, and deterministic serialization.  A degenerate
Schmidt block must be represented by a stable subspace identity.

\subsection{Benchmark TN-D: H1 variational upper bound}

For every small C8/H1 Hamiltonian block, compare exact, Krylov, exact-state TTN,
and variational TTN solutions.  Verify the upper bound, nested-bond energy
monotonicity, residual convergence, and multistart stability.

\subsection{Benchmark TN-E: observable-sensitive compression}

Choose a state for which a low-weight \(L_z\neq0\) block has a small effect on
energy but a leading effect on an OAM-sensitive current.  An energy-only bond
policy must miss the observable, while the required-sector or observable-
weighted policy retains it.

\subsection{Benchmark TN-F: common operator environment}

From one state network, evaluate charge, a current at nonzero transfer, a T-even
GTMD component, and an OAM moment.  Compare shared-environment and independent
contractions.  All values and microscopic-member identities must agree.

\subsection{Benchmark TN-G: common-parent reduction closure}

Construct a finite off-forward parent and verify the TMD, regulated GPD, PDF,
current, and Wigner routes using the same state, recoil, operator, and
normalization.  Inject a rank, mass, Fourier-phase, or momentum-fiber mismatch
and require rejection before contraction.

\subsection{Benchmark TN-H: Wilson operator plugin}

Attach the C5/C6 validation Wilson network to a compatible analytic state TTN.
Verify zero-coupling, no-cut, no-OAM, future/past, ordered-link, \(f/d\), and
soft-overlap limits.  The state topology must remain unchanged and the Wilson
path must remain operator metadata.

\subsection{Benchmark TN-I: provenance two-cell normalization}

Create two execution paths representing an explicit sector and its induced
operator, together with a matching remainder.  Verify that the compiler emits
one representative plus the remainder, that removing the two-cell produces a
double-counting failure, and that a genuinely distinct interference path is
not removed.

\subsection{Benchmark TN-J: assumption compiler}

Compile the H1 branches:
\begin{equation}
 \begin{array}{ll}
 \text{A}:&\text{induced confinement + color-spin},\\
 \text{B}:&\text{zero confinement + color-spin},\\
 \text{C}:&\text{induced confinement + no color-spin}.
 \end{array}
\label{eq:h1-branches-benchmark}
\end{equation}
Each plan must have a distinct identity and share only legitimately common
nodes.  Attempting to add A and B as simultaneous Hamiltonian mechanisms must
fail.

\subsection{Benchmark TN-K: differentiable contraction}

Compare automatic, adjoint, finite-difference, and analytic derivatives for a
small tensor contraction and a nondegenerate eigenstate.  Test a singular-value
crossing and require a subspace derivative or unavailable status rather than an
unstable vector derivative.

\subsection{Benchmark TN-L: conditional uncertainty}

Generate two assumption branches with known within-branch covariance and
between-branch displacement.  Verify \cref{eq:total-covariance-assumptions},
while also demonstrating the default report that keeps branch predictions
separate when no averaging rule is declared.

\subsection{Benchmark TN-M: contraction-plan invariance}

Contract one irregular operator network using several exact paths and slicing
schemes.  Verify identical values, explicit cost differences, deterministic
fermion signs, and agreement with a dense oracle.  An approximate route must
return a separate error record.

\section{Mandatory negative tests}
\label{sec:negative-tests}

At minimum, the implementation must reject or detect the following failures.
Stable identifiers should be assigned in the machine-readable manifest.

\begin{longtable}{L{0.10\textwidth}L{0.82\textwidth}}
\toprule
ID & Deliberate failure\\
\midrule
\endhead
N01 & Use a state tensor as a provenance relation.\\
N02 & Interpret a virtual tensor gauge transform as a QCD gauge transform.\\
N03 & Treat a contraction path as a Wilson path.\\
N04 & Mutate a frozen assumption bundle in place.\\
N05 & Reuse an assumption identity after changing one field.\\
N06 & Treat sibling physical branches as nested refinements.\\
N07 & Add two exclusive confinement branches to one Hamiltonian.\\
N08 & Select an explicit Fock sector together with its induced replacement.\\
N09 & Add a benchmark-only mechanism to a central plan.\\
N10 & Compile evolution without LF-to-QCD matching.\\
N11 & Compile a physical process with broken or unassessed factorization.\\
N12 & Promote a partial plan to validated status.\\
N13 & Add a missing physical sector through dependency closure.\\
N14 & Return an unstructured compiler error without rule identity.\\
N15 & Connect edges with matching dimensions but different physical types.\\
N16 & Omit an outer multiplicity label in a repeated irrep channel.\\
N17 & Store forbidden symmetry blocks densely and call them symmetry preserving.\\
N18 & Use a bosonic swap where a fermionic crossing is required.\\
N19 & Apply the fundamental generator to an antiquark edge.\\
N20 & Construct a color-singlet \(qqq\) times a free gluon.\\
N21 & Merge the two \(qqqg\) singlet multiplicities.\\
N22 & Remove one of the three \(qqqq\bar q\) singlet channels.\\
N23 & Change topology without updating leaf order and swap ledger.\\
N24 & Compare finite-rank topologies as though exactly equivalent.\\
N25 & Claim exact tensorization with a bond smaller than the Schmidt rank.\\
N26 & Delete the only symmetry sector required by a requested observable.\\
N27 & Call SVD compression a variational solver without Rayleigh optimization.\\
N28 & Return a variational energy below the exact finite-space ground state.\\
N29 & Violate nested-bond energy monotonicity without a convergence warning.\\
N30 & Converge energy while current or OAM errors remain above tolerance and declare completion.\\
N31 & Change an eigenstate label by eigenvalue ordering through an avoided crossing.\\
N32 & Canonicalize a degenerate Schmidt block with nondeterministic vectors and no subspace identity.\\
N33 & Reuse a contraction environment across different microscopic members.\\
N34 & Reuse a contraction environment across different operator schemes.\\
N35 & Apply a forward operator to incompatible off-forward fibers.\\
N36 & Reimplement the recoil map inside one projector.\\
N37 & Apply a scalar \(J_0\) transform to a rank-one or rank-two result.\\
N38 & Use \(\bD\) as \(\bM\) or vice versa.\\
N39 & Introduce a named-TMD-specific tensor parameter with no operator owner.\\
N40 & Interpret a missing amplitude block as a physical zero.\\
N41 & Apply the C5/C6 Wilson sign by process name rather than path identity.\\
N42 & Reverse only one member of an ordered gluon link pair when full reversal is required.\\
N43 & Mix \(f\)- and \(d\)-type gluon channels without a process color map.\\
N44 & Count the same cut in both eikonal and resolvent paths without a two-cell.\\
N45 & Omit the soft-overlap subtraction.\\
N46 & Apply the soft-overlap subtraction twice.\\
N47 & Treat provenance homology as a numerical amplitude.\\
N48 & Remove a genuinely distinct interference path using an unrelated count-once cell.\\
N49 & Differentiate a degenerate eigenvector as though nondegenerate.\\
N50 & Differentiate through an SVD crossing without a validity record.\\
N51 & Propagate raw tensor-parameter covariance including pure gauge directions.\\
N52 & Merge bond, basis, Fock, and Hamiltonian uncertainties into one unlabeled band.\\
N53 & Average mutually exclusive branches without a declared model-comparison rule.\\
N54 & Add separate branch components rather than average complete predictions.\\
N55 & Use a cached tensor whose source hash differs.\\
N56 & Reuse a cache entry after changing numerical precision or tolerance.\\
N57 & Lose deterministic reduction order in a distributed contraction.\\
N58 & Report an approximate contraction without an error estimator.\\
N59 & Connect a valence-only state to antiquark or gluon prediction readiness.\\
N60 & Connect a nucleon state to nuclear readiness without correlated proton/neutron members and nuclear operators.\\
N61 & Connect an unmatched regulated GTMD to TMD evolution.\\
N62 & Attach a process hard factor from a different scheme.\\
N63 & Modify one of the accepted 216 production reduction routes from the tensor-network package.\\
N64 & Modify an authoritative production artifact or pinned historical manifest.\\
N65 & Treat an oscillator scale as a TMD nonperturbative width.\\
N66 & Treat bond dimension as a physical parameter to be fitted to data without convergence control.\\
N67 & Insert a generic neural tensor with indices lacking physical interpretation.\\
N68 & Use a quantum circuit with no exact map to the typed Fock basis.\\
N69 & Promote a validation-only result to a physical prediction.\\
N70 & Omit assumption identity from a released TMD table.\\
N71 & Omit the compiler certificate from a released prediction plan.\\
N72 & Claim completion from test count alone while a formal acceptance gate remains unresolved.\\
\bottomrule
\end{longtable}

\section{Staged implementation program}
\label{sec:staged-program}

\subsection{P0: compiler spine during H1}

While C8/H1 is being implemented:
\begin{enumerate}
 \item implement frozen \texttt{H1AssumptionBundle} and
 \texttt{H1PredictionPlan};
 \item implement the general \AssumptionBundle, choice-axis, rule, certificate,
 and prediction-plan interfaces without connecting them to production;
 \item tensorize the exact H1 benchmark eigenvectors on one or more symmetry-
 adapted trees;
 \item establish full-rank reconstruction, canonical gauge, and serialization;
 \item construct the first provenance subcomplex for the H1 branches.
\end{enumerate}
The output remains valence-only and validation-only.

\subsection{P1: variational H1 tensor network}

After exact H1 diagonalization is stable:
\begin{enumerate}
 \item implement local TTN effective Hamiltonians and sweeps;
 \item verify upper bounds and nested-bond monotonicity;
 \item add observable-sensitive bond policies;
 \item compare exact, Krylov, and TTN states and currents;
 \item export the first tensor-network solver certificate.
\end{enumerate}
This stage validates the state engine, not a physical nucleon.

\subsection{P2: H2 dynamical gluon and sector network}

Extend the state network to
\begin{equation}
 \Hil_{qqq}\oplus\Hil_{qqqg}
\label{eq:p2-fock-space}
\end{equation}
with canonical quark--gluon vertices, instantaneous partners, sector-dependent
renormalization, complete color multiplicities, and a shared current.  The
network topology must permit intersector amplitude flow without hiding it in a
single effective valence tensor.

\subsection{P3: H3 light sea and chiral network}

Add \(qqqu\bar u\) and \(qqqd\bar d\) sectors, complete antisymmetry, chiral
interactions, and current/PCAC constraints.  Compare explicit sectors with
induced operators through provenance two-cells.  The assumption compiler must
support explicit-sector and induced-sector branches without simultaneous
counting.

\subsection{P4: microscopic common GTMD contraction}

Attach the Volume~II operator bundle to the microscopic tensor state:
\begin{enumerate}
 \item nonzero-transfer quark, antiquark, and gluon matrix elements;
 \item complete helicity matrices;
 \item common TMD/GPD/PDF/current reductions;
 \item OAM and spin--orbit moments;
 \item tensor-network, basis, Fock, and regulator convergence.
\end{enumerate}
This is the first stage that may earn \texttt{GTMD\_OVERLAP\_READY}.

\subsection{P5: Wilson, nuclear, matching, and inference plugins}

Connect the validated microscopic state to:
\begin{enumerate}
 \item the C5/C6 Wilson and cut engine, extended to higher order;
 \item the matched spin-1 nuclear state of Volume~IV;
 \item the LF-to-QCD and evolution machinery of Volume~V;
 \item the shared posterior and withheld-validation system of Volume~VI.
\end{enumerate}
Every plugin connection requires a fresh compilation certificate and
provenance complex.

\subsection{P6: distributed prediction service}

Only after the previous stages close, build a distributed service that accepts
assumption and target requests and returns conditional prediction ensembles.
The service must expose the plan, source hashes, residuals, readiness, and
uncertainty decomposition together with the numerical results.

\section{Risk register and safeguards}
\label{sec:risks}

\begin{longtable}{L{0.22\textwidth}L{0.34\textwidth}L{0.35\textwidth}}
\toprule
Risk & Failure mode & Required safeguard\\
\midrule
\endhead
Decorative tensor networks & Dense arrays are redrawn as diagrams without symmetry blocks, variational compression, or convergence. & Require typed block tensors, exact reconstruction, a true Rayleigh solver, and bond manifests.\\
Low-rank overconfidence & Energy converges while OAM, color, sea, or GTMD-sensitive blocks are lost. & Observable-sensitive bond floors and independent convergence for requested matrix elements.\\
Topology bias & One fusion tree suppresses correlations favored by another. & Full-rank equivalence tests, finite-rank topology variants, and explicit topology uncertainty.\\
Gauge conflation & Virtual tensor gauge is mistaken for QCD gauge or Wilson-path freedom. & Separate types, APIs, tests, and documentation; no cross-casting.\\
Multiplicity loss & Repeated color or spin irreps are merged. & Exact outer-multiplicity indices and C7 color-count benchmarks.\\
Fermion-sign corruption & Contraction optimizer reorders legs without graded swaps. & Explicit swap ledger and dense-oracle tests.\\
Compiler overreach & Dependency closure silently adds missing physical dynamics. & Closure may add only required implementation objects; missing physics yields partial or rejected plans.\\
Branch laundering & Mutually exclusive model choices are added as uncertainty components. & Choice-axis semantics and plan-level exclusion rules.\\
Named-function migration & A failed TMD receives its own tensor parameter. & Parameter ownership tied to Hamiltonian, current, Wilson, matching, or discrepancy operators.\\
Provenance formalism without enforcement & Two-cells are documented but not used by composition. & Executable graph normalization and deliberate double-counting injections.\\
Homology overstatement & Audit cycles are described as topological QCD observables. & Limited-topology principle and no amplitude map from homology.\\
Differentiation instability & Gradients cross singular-value or eigenvalue degeneracies. & Subspace derivatives, validity domains, and unavailable status.\\
Cache contamination & Equal-shape tensors from different schemes or members are reused. & Full content-addressed keys including source and assumption hashes.\\
Distributed nondeterminism & Parallel reductions change results or phases. & Deterministic ordering, reproducible seeds, and serial benchmark oracles.\\
Compiler/test circularity & The same implementation generates both result and expected benchmark. & Independent analytic, dense, exact-diagonalization, and alternate-contraction oracles.\\
Quantum-backend substitution & A generic circuit is called the microscopic state. & Exact basis and Hamiltonian map plus classical small-space validation.\\
Production contamination & Validation networks enter the accepted phenomenological boundary. & Disjoint provenance roots, immutable artifact hashes, and fail-closed production gates.\\
\bottomrule
\end{longtable}

\section{Final acceptance criteria}
\label{sec:acceptance}

The tensor-network and prediction-compiler layer is accepted at its first full
microscopic scope only when all of the following hold:
\begin{enumerate}
 \item Every run begins from one frozen, content-addressed assumption bundle.
 \item Compilation returns a replayable certificate and a typed prediction
 plan or a structured failure.
 \item The state, operator, and provenance networks remain distinct.
 \item Exact symmetries are structural and all required multiplicities and
 fermion signs close.
 \item Full-rank TTN tensorization reproduces exact finite-space states and a
 nontrivial operator family.
 \item The variational TTN solver satisfies upper-bound and nested-bond tests.
 \item Bond convergence is shown for energy, residual, current, OAM, and
 requested GTMD-sensitive quantities.
 \item At least two network topologies are compared at full and truncated rank.
 \item The compiler supports exclusive and nested assumption axes without
 ambiguity.
 \item Explicit-sector and induced-operator alternatives are enforced by
 executable provenance cells.
 \item All released predictions contain assumption, operator, network,
 microscopic-member, scheme, and readiness identities.
 \item One microscopic state network produces quark, antiquark, and gluon
 parents without named-function-specific parameters.
 \item TMD, GPD, PDF, current, and Wigner reductions close from the common
 parent within declared matching and numerical tolerances.
 \item Link-odd contractions use explicit Wilson operator networks and retain
 path and color identities.
 \item Nuclear, matching, evolution, process, and inference plugins remain
 fail-closed until their independent gates pass.
 \item Automatic or adjoint derivatives agree with independent benchmarks and
 carry degeneracy validity records.
 \item Tensor-network, basis, Fock, physical, matching, and statistical
 uncertainties are reported separately.
 \item At least one withheld observable family is predicted from shared
 microscopic parameters and network members.
 \item The complete negative-test suite passes, including production-isolation
 and immutable-artifact gates.
 \item Removing all present phenomenological named-TMD inputs does not make the
 microscopic tensor-network calculation undefined.
\end{enumerate}

\section{Recommended implementation boundary after C8}
\label{sec:next-boundary}

If C8/H1 completes successfully, the next package should be split into two
coordinated but independently gated work streams.

\subsection{C9A: general compiler and tensor-network core}

Implement:
\begin{itemize}
 \item general \AssumptionBundle, \texttt{ChoiceAxis}, and
 \texttt{AdmissibilityRule};
 \item general \PredictionPlan\ and compilation certificate;
 \item \TensorEdgeType, \SymTensor, \FusionTree, and bond manifests;
 \item exact H1 tensorization and canonicalization;
 \item variational TTN sweeps and observable-sensitive compression;
 \item H1 branch provenance and conditional result envelopes.
\end{itemize}
This package remains valence-only and disconnected from production TMDs.

\subsection{C9B/H2: dynamical \texorpdfstring{\(qqqg\)}{qqqg} sector}

Implement the first interacting sector extension with canonical quark--gluon
vertices, instantaneous partners, sector-dependent renormalization, current and
Ward closure, gluon/OAM exports, and a microscopic state interface compatible
with C5/C6.  C9B must consume the C9A tensor/compiler interfaces rather than
create another state representation.

\subsection{Why the split is useful}

The compiler and TTN engine are reusable mathematical infrastructure, while H2
introduces new physical dynamics and counterterms.  Separate commits and
acceptance manifests make it possible to diagnose whether a failure arises in
the tensor representation or in the Hamiltonian model.

\section{Conclusion}
\label{sec:conclusion}

The geometric and topological proposal becomes an executable predictive model
only when the model choices, quantum state, operator contractions, provenance,
and uncertainty are compiled together.  This volume supplies that missing
layer.

The central object is not a tensor network by itself.  It is the chain
\begin{equation}
 \boxed{
 \Aset
 \longrightarrow
 \Pplan
 \longrightarrow
 \Nstate
 \longrightarrow
 \{\Nop\}
 \longrightarrow
 \{W_{a/h},F_A,\cO_{\mathrm{process}}\}
 }
\label{eq:final-chain}
\end{equation}
with an accompanying provenance complex and validation certificate.  The state
network carries the microscopic amplitudes; the operator network carries the
Hamiltonian, currents, GTMDs, Wilson lines, and process insertions; the
provenance complex carries assumption and no-double-counting relations.  Their
separation prevents elegant notation from replacing dynamics, while their
typed composition makes the consequences of each declared model choice
computable and auditable.

Symmetry-adapted tensors ensure that color, flavor, spin, OAM, statistics, and
Fock identity survive compression.  Variational TTNs make bond rank an explicit
numerical truncation with upper bounds and observable-level convergence.
The assumption compiler turns alternative Hamiltonians, sectors, Wilson paths,
nuclear models, matching schemes, and process qualifications into distinct
prediction plans rather than an undifferentiated sum.  The two-complex records
where different computational paths are equivalent, subtractive, replacing,
or genuinely distinct.  Conditional ensembles then state precisely which TMD
predictions follow from which assumptions.

The immediate task is not to claim a complete tensor-network solution of QCD.
It is to validate this execution layer on the C8/H1 interacting valence state,
then enlarge the same state and compiler through \(qqqg\), light-sea, GTMD,
Wilson, nuclear, matching, evolution, and inference stages.  When those layers
close, the project will realize the objective set at the beginning: a
geometric, compositional, tensor-network model in which scientifically
admissible assumptions generate correlated and falsifiable predictions for the
full spin-1 TMD system.

\appendix

\section{Assumption-bundle schema}
\label{app:assumption-schema}

\begin{verbatim}
{
  "assumption_bundle_id": "sha256:...",
  "scope": "validation|microscopic|physical",
  "physics": {
    "hamiltonian_route": "RSA_BLFQ_EFT",
    "fock_sectors": ["qqq", "qqqg", "qqquubar", "qqqddbar"],
    "interaction_terms": ["..."],
    "current_basis": "...",
    "wilson_order": 0,
    "nuclear_route": "MATCHED_HADRONIC_EFT"
  },
  "regulators": {
    "K": "...", "Nmax": "...", "b_GeV": "...",
    "lambda_GeV": "...", "uv": "...", "ir": "...",
    "endpoint_policy": "...", "rapidity_regulator": "..."
  },
  "algebra": {
    "symmetry_contract": "...",
    "color_fusion_tree": "...",
    "fermion_order": "...",
    "tensor_topology": "...",
    "operator_basis_version": "..."
  },
  "numerics": {
    "solver": ["EXACT", "KRYLOV", "TTN"],
    "bond_policy": "...",
    "contraction_policy": "...",
    "precision": "...", "tolerances": {"...": "..."}
  },
  "process": {
    "target_observables": ["..."],
    "factorization_status": "...",
    "scheme": "...", "accuracy_manifest": "..."
  },
  "inference": {
    "calibration_partition": "...",
    "holdout_partition": "...",
    "posterior_member_set": "...",
    "model_averaging_rule": null
  },
  "source_hashes": {"...": "sha256:..."},
  "schema_version": "V8.1"
}
\end{verbatim}

\section{Prediction-plan schema}
\label{app:plan-schema}

\begin{verbatim}
{
  "prediction_plan_id": "sha256:...",
  "assumption_bundle_id": "sha256:...",
  "target_id": "...",
  "compilation_certificate_id": "...",
  "nodes": [
    {
      "node_id": "...",
      "type": "basis|hamiltonian|solver|state|operator|reduction|result",
      "inputs": ["..."],
      "object_id": "...",
      "provenance_node_id": "...",
      "readiness_required": ["..."],
      "cache_key": "sha256:..."
    }
  ],
  "edges": [{"source": "...", "target": "...", "map_id": "..."}],
  "validation_gates": ["..."],
  "holdout_observables": ["..."],
  "execution_manifest_id": "...",
  "status": "REJECTED|PARTIAL|EXECUTABLE|VALIDATED"
}
\end{verbatim}

\section{Symmetry-block tensor schema}
\label{app:tensor-schema}

\begin{verbatim}
{
  "tensor_id": "sha256:...",
  "role": "state|operator|isometry|environment",
  "edge_types": ["..."],
  "fusion_tree_id": "...",
  "canonical_gauge": "...",
  "blocks": [
    {
      "charge_tuple": ["..."],
      "outer_multiplicity": "...",
      "structural_tensor_id": "...",
      "degeneracy_shape": ["..."],
      "data_hash": "sha256:..."
    }
  ],
  "fermion_swap_ledger_id": "...",
  "resolution_id": "...",
  "microscopic_member_id": "...",
  "schema_version": "V8.1"
}
\end{verbatim}

\section{Bond and solver manifest}
\label{app:bond-schema}

\begin{verbatim}
{
  "state_tensor_network_id": "...",
  "topology_id": "...",
  "orthogonality_center": "...",
  "edges": [
    {
      "edge_id": "...",
      "sector_dimensions": {"charge": "chi_degeneracy"},
      "required_floors": {"charge": "minimum"},
      "schmidt_values_hash": "...",
      "discarded_weight": "...",
      "entropy": "..."
    }
  ],
  "solver": {
    "route": "EXACT|KRYLOV|VARIATIONAL_TTN",
    "energy": "...",
    "residual": "...",
    "sweeps": "...",
    "initialization": "...",
    "multistart_manifest": "..."
  },
  "observable_errors": {"...": "..."},
  "exact_oracle_id": "..."
}
\end{verbatim}

\section{Provenance two-complex schema}
\label{app:provenance-schema}

\begin{verbatim}
{
  "complex_id": "sha256:...",
  "assumption_bundle_id": "...",
  "zero_cells": [{"id": "...", "type": "state|operator|scheme|result"}],
  "one_cells": [
    {"id": "...", "source": "...", "target": "...", "relation": "..."}
  ],
  "two_cells": [
    {
      "id": "...",
      "boundary_paths": [["..."], ["..."]],
      "relation": "EQUIVALENT_COUNT_ONCE|SUBTRACT|REPLACES|ALTERNATIVE",
      "coefficient": "...",
      "remainder_path": ["..."],
      "validity_domain": "...",
      "tolerance": "..."
    }
  ],
  "normal_form_rules": ["..."],
  "unresolved_homology_classes": ["..."]
}
\end{verbatim}

\section{Compact adoption checklist}
\label{app:checklist}

Before calling a result an assumption-conditioned tensor-network prediction,
answer:
\begin{enumerate}
 \item Is the complete assumption bundle frozen and recoverable?
 \item Did the compiler pass every dependency, exclusion, scheme, and readiness
 rule?
 \item Are state, operator, and provenance networks separate?
 \item Does every edge have a physical or declared numerical meaning?
 \item Are exact symmetries and fermionic signs structural?
 \item Are all outer multiplicities present?
 \item Does full-rank tensorization reproduce the exact finite state?
 \item Was a genuine variational TTN optimization performed?
 \item Do bond and topology convergence include nonenergy observables?
 \item Are Wilson paths and tensor virtual gauges distinguished?
 \item Are all GTMD/TMD reductions traced to one common state and operator
 bundle?
 \item Are explicit and induced descriptions prevented from double counting?
 \item Does every derivative carry a validity record?
 \item Are within- and between-assumption uncertainties separated?
 \item Is the released result accompanied by its compiler and solver
 certificates?
 \item Are unsupported downstream statuses fail-closed?
 \item Can the calculation be replayed from immutable source and content hashes?
 \item Does at least one withheld observable test the shared microscopic
 prediction?
\end{enumerate}

\small
\begin{thebibliography}{99}
\raggedright

\bibitem{Penrose1971}
R.~Penrose,
``Applications of negative dimensional tensors,''
in \emph{Combinatorial Mathematics and its Applications},
Academic Press, New York (1971), pp.~221--244.

\bibitem{White1992}
S.~R.~White,
``Density matrix formulation for quantum renormalization groups,''
Phys. Rev. Lett. \textbf{69}, 2863 (1992).

\bibitem{Shi2006}
Y.-Y.~Shi, L.-M.~Duan, and G.~Vidal,
``Classical simulation of quantum many-body systems with a tree tensor
network,''
Phys. Rev. A \textbf{74}, 022320 (2006),
arXiv:quant-ph/0511070.

\bibitem{Verstraete2008}
F.~Verstraete, V.~Murg, and J.~I.~Cirac,
``Matrix product states, projected entangled pair states, and variational
renormalization group methods for quantum spin systems,''
Adv. Phys. \textbf{57}, 143 (2008),
arXiv:0907.2796.

\bibitem{Orus2014}
R.~Or\'us,
``A practical introduction to tensor networks: Matrix product states and
projected entangled pair states,''
Ann. Phys. \textbf{349}, 117 (2014),
arXiv:1306.2164.

\bibitem{Singh2009}
S.~Singh, R.~N.~C.~Pfeifer, and G.~Vidal,
``Tensor network decompositions in the presence of a global symmetry,''
Phys. Rev. A \textbf{82}, 050301(R) (2010),
arXiv:0907.2994.

\bibitem{Singh2011}
S.~Singh, R.~N.~C.~Pfeifer, and G.~Vidal,
``Tensor network states and algorithms in the presence of a global
\(U(1)\) symmetry,''
Phys. Rev. B \textbf{83}, 115125 (2011),
arXiv:1008.4774.

\bibitem{Weichselbaum2012}
A.~Weichselbaum,
``Non-Abelian symmetries in tensor networks: A quantum symmetry space
approach,''
Ann. Phys. \textbf{327}, 2972 (2012),
arXiv:1202.5664.

\bibitem{Weichselbaum2020}
A.~Weichselbaum,
``X-symbols for non-Abelian symmetries in tensor networks,''
Phys. Rev. Research \textbf{2}, 023385 (2020),
arXiv:1910.13736.

\bibitem{Gunst2018}
K.~Gunst, F.~Verstraete, S.~Wouters, O.~Legeza, and D.~Van Neck,
``T3NS: Three-legged tree tensor network states,''
J. Chem. Theory Comput. \textbf{14}, 2026 (2018),
arXiv:1801.09998.

\bibitem{Haegeman2014}
J.~Haegeman, M.~Mari\"en, T.~J.~Osborne, and F.~Verstraete,
``Geometry of matrix product states: Metric, parallel transport, and
curvature,''
J. Math. Phys. \textbf{55}, 021902 (2014),
arXiv:1210.7710.

\bibitem{Haegeman2013}
J.~Haegeman, T.~J.~Osborne, and F.~Verstraete,
``Post-matrix product state methods: To tangent space and beyond,''
Phys. Rev. B \textbf{88}, 075133 (2013),
arXiv:1305.1894.

\bibitem{Liao2019}
H.-J.~Liao, J.-G.~Liu, L.~Wang, and T.~Xiang,
``Differentiable programming tensor networks,''
Phys. Rev. X \textbf{9}, 031041 (2019),
arXiv:1903.09650.

\bibitem{Tagliacozzo2014}
L.~Tagliacozzo, A.~Celi, and M.~Lewenstein,
``Tensor networks for lattice gauge theories with continuous groups,''
Phys. Rev. X \textbf{4}, 041024 (2014),
arXiv:1405.4811.

\bibitem{Canals2024}
M.~Canals, N.~Chepiga, and L.~Tagliacozzo,
``A tensor network formulation of lattice gauge theories based only on
symmetric tensors,''
arXiv:2412.16961 (2024).

\bibitem{GrayKourtis2021}
J.~Gray and S.~Kourtis,
``Hyper-optimized tensor network contraction,''
Quantum \textbf{5}, 410 (2021),
arXiv:2002.01935.

\bibitem{Acuaviva2023}
A.~Acuaviva, V.~Makam, H.~Nieuwboer, D.~P\'erez-Garc\'ia,
F.~Sittner, M.~Walter, and F.~Witteveen,
``The minimal canonical form of a tensor network,''
arXiv:2209.14358 (2022).

\bibitem{Fong2015}
B.~Fong,
``Decorated cospans,''
Theory Appl. Categ. \textbf{30}, 1096 (2015),
arXiv:1502.00872.

\bibitem{Courser2017}
K.~Courser,
``A bicategory of decorated cospans,''
Theory Appl. Categ. \textbf{32}, 995 (2017),
arXiv:1605.08100.

\bibitem{BaezCourser2020}
J.~C.~Baez and K.~Courser,
``Structured cospans,''
Theory Appl. Categ. \textbf{35}, 1771 (2020),
arXiv:1911.04630.

\bibitem{Courser2020}
K.~Courser,
``Open systems: A double categorical perspective,''
arXiv:2008.02394 (2020).

\bibitem{FongSpivak2019}
B.~Fong and D.~I.~Spivak,
``Hypergraph categories,''
J. Pure Appl. Algebra \textbf{223}, 4746 (2019),
arXiv:1806.08304.

\bibitem{FongSpivak2018}
B.~Fong and D.~I.~Spivak,
\emph{Seven Sketches in Compositionality: An Invitation to Applied Category
Theory},
arXiv:1803.05316 (2018).

\bibitem{Dirac1949}
P.~A.~M.~Dirac,
``Forms of relativistic dynamics,''
Rev. Mod. Phys. \textbf{21}, 392 (1949).

\bibitem{BrodskyPauliPinsky1998}
S.~J.~Brodsky, H.-C.~Pauli, and S.~S.~Pinsky,
``Quantum chromodynamics and other field theories on the light cone,''
Phys. Rept. \textbf{301}, 299 (1998),
arXiv:hep-ph/9705477.

\bibitem{Vary2010}
J.~P.~Vary et al.,
``Hamiltonian light-front field theory in a basis function approach,''
Phys. Rev. C \textbf{81}, 035205 (2010),
arXiv:0905.1411.

\bibitem{Xu2025}
S.~Xu, Y.~Liu, C.~Mondal, J.~Lan, X.~Zhao, Y.~Li, and J.~P.~Vary,
``Towards a first principles light-front Hamiltonian for the nucleon,''
Phys. Lett. B \textbf{867}, 139599 (2025),
arXiv:2408.11298.

\bibitem{Karmanov2008}
V.~A.~Karmanov, J.-F.~Mathiot, and A.~V.~Smirnov,
``Systematic renormalization scheme in light-front dynamics with Fock space
truncation,''
Phys. Rev. D \textbf{77}, 085028 (2008),
arXiv:0801.4507.

\bibitem{BrodskyDiehlHwang2001}
S.~J.~Brodsky, M.~Diehl, and D.~S.~Hwang,
``Light-cone wavefunction representation of deeply virtual Compton scattering,''
Nucl. Phys. B \textbf{596}, 99 (2001),
arXiv:hep-ph/0009254.

\bibitem{Meissner2009}
S.~Mei{\ss}ner, A.~Metz, and M.~Schlegel,
``Generalized parton correlation functions for a spin-1/2 hadron,''
JHEP \textbf{08}, 056 (2009),
arXiv:0906.5323.

\bibitem{Lorce2011}
C.~Lorc\'e and B.~Pasquini,
``Quark Wigner distributions and orbital angular momentum,''
Phys. Rev. D \textbf{84}, 014015 (2011),
arXiv:1106.0139.

\bibitem{Collins2011}
J.~Collins,
\emph{Foundations of Perturbative QCD},
Cambridge University Press (2011).

\end{thebibliography}

\end{document}
